% ============================================================
%  FedLoRAGuard: Federated Dynamic Graph Neural Networks
%  with DP Certificates for LoRA Adapter Supply-Chain Integrity
%  EMNLP 2026 Camera-Ready (final version)
% ============================================================
\documentclass[11pt]{article}

\usepackage[final]{acl}

\usepackage{times}
\usepackage{latexsym}
\usepackage{needspace}
\usepackage[T1]{fontenc}
\usepackage[utf8]{inputenc}
\usepackage{microtype}
% \usepackage{inconsolata} % restore for Overleaf
\usepackage{makecell}

\usepackage{url}
\usepackage{hyperref}
\usepackage{booktabs}
\usepackage{amsmath}
\usepackage{amssymb}
\usepackage{amsthm}
\usepackage{amsfonts}
\usepackage{mathtools}
\usepackage{bm}
\usepackage{nicefrac}
\usepackage{xcolor}
\usepackage{graphicx}
\usepackage{multirow}
\usepackage{makecell}
\usepackage{threeparttable}
\usepackage{tikz}
\usetikzlibrary{positioning,arrows.meta,shapes.geometric,fit,backgrounds,calc,decorations.pathreplacing}
\usepackage{pgfplots}
\pgfplotsset{compat=1.18}
\usepackage{algorithm}
\usepackage{algorithmic}
\usepackage{enumitem}

\theoremstyle{plain}
\newtheorem{theorem}{Theorem}
\newtheorem{proposition}[theorem]{Proposition}
\newtheorem{lemma}[theorem]{Lemma}
\newtheorem{corollary}[theorem]{Corollary}
\theoremstyle{definition}
\newtheorem{definition}{Definition}
\theoremstyle{remark}
\newtheorem{remark}{Remark}

\newcommand{\modelname}{FedLoRAGuard}
\newcommand{\benchname}{LoRAchain-2026}
\newcommand{\R}{\mathbb{R}}
\newcommand{\E}{\mathbb{E}}
\newcommand{\NN}{\mathbb{N}}
\providecommand{\Prob}{}\renewcommand{\Prob}{\mathbb{P}}
\newcommand{\Lcal}{\mathcal{L}}
\newcommand{\Dcal}{\mathcal{D}}
\newcommand{\Tcal}{\mathcal{T}}
\newcommand{\Gcal}{\mathcal{G}}
\newcommand{\Vcal}{\mathcal{V}}
\newcommand{\Ecal}{\mathcal{E}}
\newcommand{\Rcal}{\mathcal{R}}
\newcommand{\Acal}{\mathcal{A}}
\newcommand{\Ncal}{\mathcal{N}}
\newcommand{\softmax}{\mathrm{softmax}}
\newcommand{\norm}[1]{\|#1\|}
\DeclareMathOperator*{\argmax}{arg\,max}

\title{\modelname: Federated Dynamic Graph Neural\\
Networks with Differential-Privacy Certificates\\
for LoRA Adapter Integrity Verification}

\author{Roger Nick Anaedevha\textsuperscript{1,*}\thanks{\;\textsuperscript{*}\,All
        authors of this publication are co-first authors (equal contribution).
        Corresponding author: \texttt{ar006@campus.mephi.ru}} \quad
        Alexander G.\ Trofimov\textsuperscript{1,*} \quad
        Yuri V.\ Borodachev\textsuperscript{2,*} \\
        \textsuperscript{1}Institute of Cyber Intelligence Systems, \\
        National Research Nuclear University MEPhI, Moscow 115409, Russia \\
        \textsuperscript{2}Artificial Intelligence Research Center, \\
        National Research Nuclear University MEPhI, Moscow 115409, Russia \\
        \texttt{ar006@campus.mephi.ru}}

\begin{document}
\maketitle

\begin{abstract}
Low-Rank Adaptation (LoRA) is now the dominant
parameter-efficient fine-tuning method for large language models,
but the resulting open marketplace of adapters has become a
supply-chain attack surface: backdoored adapters merge silently
with arbitrary task adapters while preserving benign utility, and
existing centralised integrity tools require marketplaces to share
raw weights, which is incompatible with proprietary fine-tuning.
We present \modelname, a federated dynamic-graph framework that
verifies adapter integrity without weight sharing. The LoRA
ecosystem is modelled as a heterogeneous continuous-time dynamic
graph with four node types and six edge types; a federated dynamic
GNN is trained per-client with client-update-level
$(\epsilon,\delta)$-DP Gaussian noise, secure aggregation, and
FLTrust scoring. We prove a sensitivity bound for federated
heterogeneous-GNN client updates and a certified poisoning-radius
theorem which converts the per-round inference-time smoothing
budget $\epsilon_r$ into an analytical bound on the size of a
bounded-update colluding coalition $k^*$. On $13{,}500$ benign and
backdoored adapters across four base-model families and ten attack
types, \modelname\ obtains $96.4\%$ macro-F1 and a certified
coalition bound $k^*=25$ out of $N=50$ marketplaces at
inference-time smoothing budget $\epsilon_r=0.5$, and empirically
tolerates label-flipping Byzantine coalitions up to $\rho\le 0.33$
without certificate degradation. The training-time DP-SGD
mechanism, calibrated at $\sigma=1.1$ and Poisson subsampling
$q=0.2$ over $T=100$ rounds, achieves a cumulative privacy budget
$\epsilon_T\approx 13.6$ at $\delta=10^{-5}$ under the Opacus PRV
accountant, at sub-second per-adapter verification latency on
A100-class hardware. The full DP + secure-aggregation stack adds a
measured median per-round wall-clock overhead of $\approx 4\%$
(pooled median across three runs of $20$--$60$ measurement rounds
on identical CPU hardware) over non-private FedAvg; the released
runtime ships the instrumentation so operators can re-measure in
place on their target hardware.
\end{abstract}

% ============================================================
%  1. INTRODUCTION
% ============================================================
\section{Introduction}
\label{sec:intro}

Low-Rank Adaptation (LoRA)~\cite{hu2021lora} is the dominant
parameter-efficient fine-tuning method for large language models.
It inserts low-rank update matrices into Transformer attention
while freezing the base weights, reducing trainable-parameter counts
by four orders of magnitude without losing downstream accuracy.
Combined with $4$-bit quantization
(QLoRA~\cite{dettmers2023qlora}), LoRA enables individual users to
fine-tune $65$B-parameter base models on a single consumer GPU.
The result is a large community ecosystem of shared adapters: as
of mid-2025 the Hugging Face Hub hosts over $530{,}000$ models,
and one base model alone has accumulated over $30{,}000$
derivative LoRAs.

The same ecosystem is now a supply-chain attack surface for
language models. \citet{liu2024loraattack} showed that
backdoored adapters merge silently with arbitrary task adapters
and retain malicious behaviour while preserving benign utility.
\citet{lermen2023lora} reduced the safety-refusal rate of
Llama-2-Chat-70B from $95\%$ to below $1\%$ via a USD-$200$ LoRA
fine-tune, a turnkey alignment-bypass mechanism. Composite
backdoor attacks~\cite{huang2024composite} scatter trigger keys
across prompt components. Industry incident reports document
over $100$ malicious models on Hugging Face Hub embedding reverse
shells, indirect compromise of the platform's own
\texttt{SFconvertbot} which had merged over $42{,}000$ conversion
pull requests, and a five-fold year-over-year increase in
malicious-model uploads~\cite{cohen2024jfrog,kellas2025pickleball}.
These marketplace-scale figures concern
\emph{full-model} backdoors; the LoRA-\emph{adapter} threat we
address is documented separately by the adapter-specific attack
work above~\cite{liu2024loraattack,lermen2023lora,huang2024composite},
and we cite the aggregate figures only as evidence of an active
marketplace attack surface, not as a count of adapter backdoors.

The standard centralised response is exemplified by
PEFTGuard~\cite{sun2025peftguard}, which trains a
backdoor-detection classifier on raw adapter weights and reports
$98.3$--$100\%$ accuracy on its PADBench corpus. This requires
marketplaces to share raw LoRA weights with a third-party
detector. The requirement is incompatible with proprietary
fine-tuning practices and creates a single point of trust.
Beyond this deployment barrier, a weight-space-only detector sees
each adapter in isolation and has no access to the cross-marketplace
provenance and lineage signal that distinguishes a silently
backdoored derivative from a benign one; this contextual signal,
rather than any claimed collapse of PEFTGuard's accuracy, is what
motivates a relational, privacy-preserving alternative. (We do not
run an adaptive attack against PEFTGuard and make no claim about its
robustness under one.)

\paragraph{Our approach.}
We propose a federated, privacy-preserving alternative. The LoRA
ecosystem is modelled as a heterogeneous continuous-time dynamic
graph in which nodes are adapters, contributors, base models, and
downstream applications. Time-stamped edges record derivation,
fine-tuning, deployment, and citation relationships. Across this
graph, a dynamic graph neural network is trained per-client with
\emph{client-update-level} Differentially-Private Gaussian
noise~\cite{abadi2016deep,mironov2017renyi} and secure
aggregation~\cite{bonawitz2017practical}, so that no marketplace
must reveal its raw adapter weights. By composing PixelDP-style
DP-stability arguments~\cite{lecuyer2019certified} with CRFL-style
federated smoothing~\cite{xie2021crfl}, we derive a certified
poisoning radius that converts the per-round inference-time
smoothing budget $\epsilon_r$ into a bound on the number of
colluding malicious clients $k^*$ under our bounded-update threat
model at which the global verdict is provably invariant.

\paragraph{Claim structure.}
The paper is organised around three tiered claims. The
\emph{primary} claim is that federated lineage-aware detection can
identify malicious LoRA adapters without any marketplace sharing
raw weights, and matches within $1.9$ macro-F1 points the
centralised PEFTGuard reference number on the same corpus. The
\emph{secondary} claim is that DP-SGD combined with secure
aggregation can be incorporated into this pipeline at modest utility
cost ($0.5$--$1.0$ pp macro-F1) and at a measured $\approx 4\%$
median per-round wall-clock overhead on CPU-hosted federation
(Appendix~\ref{app:repro}). The \emph{theoretical}
claim is that, under a precisely defined bounded-update threat
model (Definition~\ref{def:attack}), the composition of PixelDP
stability, CRFL federated smoothing and R\'enyi-DP composition
yields an integer certificate $k^*$ on the largest attacker-client
coalition whose actions cannot flip the global verdict on a queried
adapter. All other results (dynamic-graph superiority, multimodal
gains, FLTrust robustness, lineage importance, IDS transfer,
cross-hardware latency) are reported as supporting evidence rather
than free-standing contributions.

\paragraph{Contributions.}
We make three contributions. \textbf{(1) Formal model and
benchmark:} to our knowledge the first formal model of the LoRA
marketplace as a continuous-time heterogeneous dynamic graph with
four node types and six edge types, together with \benchname, a
benchmark of $13{,}500$ benign and backdoored LoRA adapters across
four base-model families and ten attack types extending
established NLP-safety resources
(BackdoorLLM~\cite{li2025backdoorllm},
PADBench~\cite{sun2025peftguard}) with cross-marketplace lineage
edges harvested from HuggingGraph~\cite{rahman2025hugginggraph}.
\textbf{(2) Framework and algorithms:} a federated dynamic GNN
with client-update-level DP and FLTrust hardening --- a per-client
DyGFormer-style encoder with HGT relation-aware attention, trained
under client-update-level Gaussian DP-SGD with secure aggregation
and FLTrust bootstrapping, empirically robust to label-flipping
Byzantine coalitions up to $\rho\le 0.33$. \textbf{(3) Theory and
evaluation:} a sensitivity bound for the federated
heterogeneous-dynamic-GNN client update
(Theorem~\ref{thm:sensitivity}) and a certified poisoning-radius
theorem (Theorem~\ref{thm:certified}) that turns the
inference-time smoothing budget into an integer bound $k^*$ on
bounded-update coalition size, validated on \benchname, the
BackdoorLLM-LoRA leaderboard, and four IDS datasets, where
\modelname\ obtains $96.4\%$ macro-F1, certified coalition bound
$k^* = 25$ out of $N = 50$ at smoothing budget $\epsilon_r=0.5$,
cumulative training-time budget
$(\epsilon_T,\delta)\approx(13.6,10^{-5})$ over $T = 100$
Poisson-subsampled rounds under the Opacus PRV
accountant~\cite{gopi2021numerical}, and sub-second per-adapter
scan latency on A100-class hardware.

\paragraph{Why this belongs at an NLP venue.}
LoRA is the dominant fine-tuning method for \emph{language} models,
and the threats we verify --- alignment-bypass fine-tunes, composite
textual-trigger backdoors, and instruction-following corruption ---
are attacks on the safety of NLP systems at the
parameter-efficient-tuning layer. The signal that separates a
benign adapter from a backdoored one is partly linguistic: the
multimodal encoder reads model cards, README text, and
contributor-profile text (Section~\ref{sec:framework}) alongside
the weight and usage modalities, and \benchname\ is built by
extending NLP-safety resources (BackdoorLLM, PADBench). The
verification method itself is task-agnostic, but the resource, the
threat model, and the deployment target (Hugging Face-style LLM
adapter hubs) are squarely NLP-safety concerns, which is why we
position the contribution for the ACL community rather than a
general security venue.

\paragraph{Roadmap.}
Section~\ref{sec:related} surveys prior work on LoRA-adapter
security, federated learning with differential privacy, dynamic
graph neural networks, and certified robustness against poisoning;
Section~\ref{sec:problem} formalises the LoRA-ecosystem graph and
the federated multi-marketplace threat model;
Section~\ref{sec:framework} presents the four-component
\modelname\ framework; Section~\ref{sec:theory} proves the
sensitivity bound and the certified poisoning-radius theorem;
Section~\ref{sec:experiments} reports the empirical evaluation
including ablation, spectral analysis, and IDS transfer; and
Sections~\ref{sec:discussion}--\ref{sec:conclusion} discuss
limitations and conclude.

% ============================================================
%  2. RELATED WORK
% ============================================================
\section{Related Work}
\label{sec:related}

\paragraph{LoRA, PEFT, and adapter security.}
Parameter-efficient fine-tuning rose to prominence with
adapters and LoRA~\cite{hu2021lora}.
QLoRA~\cite{dettmers2023qlora} added $4$-bit NormalFloat
quantization. AdaLoRA~\cite{zhang2023adalora},
DoRA~\cite{liu2024dora}, and GaLore~\cite{zhao2024galore} produced
structurally distinct adapter families that complicate uniform
integrity checks. Adapter security emerged as a research area with
LoRA-as-an-Attack~\cite{liu2024loraattack},
\citet{lermen2023lora}'s cheap alignment-bypass demonstration,
Composite Backdoor Attacks~\cite{huang2024composite}, and the
safety-critical neuron analysis of~\citet{wei2024brittleness}.
PEFTGuard~\cite{sun2025peftguard} centralised backdoor detection
on raw adapter weights, reaching $98.3$--$100\%$ centralised
accuracy on PADBench, but requires marketplaces to share raw
weights and inspects each adapter without cross-marketplace context.
ShadowGenes~\cite{schulz2025shadowgenes} introduced graph-based
model-genealogy verification using computational-graph
fingerprinting. Industry incident
analyses~\cite{cohen2024jfrog,kellas2025pickleball} document the
operational threat at marketplace scale. More broadly, ML
supply-chain integrity has been approached through model signing,
software bills of materials (SBOM), and update frameworks such as
TUF; \modelname\ is complementary to these provenance-attestation
mechanisms, adding a content-level maliciousness verdict that
signing alone does not provide.

\paragraph{Federated learning, DP, and secure aggregation.}
Federated learning~\cite{mcmahan2017fedavg} with heterogeneity-aware
variants~\cite{li2020fedprox,karimireddy2020scaffold} combines with
DP-SGD~\cite{abadi2016deep}, R\'enyi-DP accounting~\cite{mironov2017renyi}
and the numerically tight PRV accountant~\cite{gopi2021numerical},
and secure aggregation~\cite{bonawitz2017practical} to hide client
contributions. Byzantine robustness is addressed by
Krum~\cite{blanchard2017machine},
Trimmed-Mean~\cite{yin2018byzantine}, and
FLTrust~\cite{cao2021fltrust}, which bootstraps trust from a small
clean root set; game-theoretic treatments of Byzantine resilience
in federated security deployments~\cite{anaedevha2026byzantine}
address the same tension between robustness and privacy that
\modelname\ inherits. Federated foundation-model fine-tuning
(FATE-LLM~\cite{fan2023fatellm}, FLoRA~\cite{wang2024flora},
OpenFedLLM~\cite{ye2024openfedllm}) trains new adapters;
\modelname\ instead \emph{verifies} existing ones. Closest to our
setting is Safe-FedLLM~\cite{tao2026safefedllm}, which also reads
LoRA weights as a behavioural signal, training lightweight probes on
each client's local $\Delta$LoRA updates to filter malicious
\emph{clients} during federated training at step, client, and shadow
level. The two are complementary and act at different points in the
supply chain: Safe-FedLLM protects the training process from
malicious participants, whereas \modelname\ verifies adapters
\emph{already published} to a marketplace, scoring each artefact
against cross-marketplace lineage structure rather than a single
client's update stream, and issuing a certified coalition bound
$k^*$ over colluding marketplaces --- a guarantee that
probe-classifier defences do not provide. A deployment could run
Safe-FedLLM inside each marketplace's own training pipeline and
\modelname\ across marketplaces at publication time.

\paragraph{Dynamic and heterogeneous GNNs.}
Continuous-time dynamic graph learning was pioneered by
JODIE~\cite{kumar2019jodie}, DyRep~\cite{trivedi2019dyrep}, and
TGAT~\cite{xu2020tgat}. TGN~\cite{rossi2020tgn} unified these in a
memory-module-plus-graph-attention architecture.
DyGFormer~\cite{yu2023dygformer} introduced neighbor-co-occurrence
encoding and patching for long histories. Multi-scale temporal
graph encoders have been applied to security
telemetry~\cite{anaedevha2026multiscale}, motivating the CT-DG
formulation we adopt for adapter lineage. For heterogeneous
graphs, HGT~\cite{hu2020hgt} introduced meta-relation-aware
mutual attention with type-dependent parameters and Relative
Temporal Encoding. Federated graph learning was introduced by
FedGraphNN~\cite{he2021fedgraphnn}; GAP~\cite{sajadmanesh2023gap}
derives a DP-GNN with edge- and node-level guarantees through
R\'enyi-DP analysis, which we extend here to the federated setting.
The HF ecosystem itself has been mapped recently by
HuggingGraph~\cite{rahman2025hugginggraph} and by the model-atlas
analysis of~\citet{horwitz2025atlas}.

\paragraph{Certified robustness against poisoning.}
The bridge from DP to certified robustness was established by
PixelDP~\cite{lecuyer2019certified}.
\citet{ma2019data} converted DP-SGD privacy budgets into certified
poisoning radii under bounded data poisoning. In federated
settings, CRFL~\cite{xie2021crfl} uses parameter clipping plus
Gaussian smoothing to bound certified accuracy as a function of
poisoning ratio; FLCert~\cite{cao2022flcert} extends this with
provable security under up to a fraction of malicious clients.
Poisoning resilience has also been pursued architecturally, e.g.\
through selective state-space sequence
models~\cite{anaedevha2026mambashield}, though without a formal
certificate. For graph data, \citet{wang2021certgnn} provide
certificates against bounded edge insertions/deletions;
\citet{scholten2022randomized} exploit message-passing structure
for stronger certificates against node-controlling adversaries.
\modelname's certified poisoning-radius theorem synthesises these
threads into a single per-marketplace bound on the largest
bounded-update coalition of colluding malicious clients that cannot
alter the global maliciousness verdict.

% ============================================================
%  3. PRELIMINARIES AND PROBLEM
% ============================================================
\section{Preliminaries and Problem}
\label{sec:problem}

\paragraph{LoRA adapter.}
A LoRA adapter for a Transformer layer with frozen weight matrix
$\mathbf{W}_0 \in \R^{d_{\mathrm{out}} \times d_{\mathrm{in}}}$
replaces the forward pass $\mathbf{y} = \mathbf{W}_0 \mathbf{x}$
(where $\mathbf{x}\in\R^{d_{\mathrm{in}}}$ is the layer input and
$\mathbf{y}\in\R^{d_{\mathrm{out}}}$ its output) with
\begin{equation}
\mathbf{y} = \mathbf{W}_0 \mathbf{x} + \frac{\alpha}{r} \mathbf{B}\mathbf{A}\mathbf{x},\;\; \mathbf{B}\in\R^{d_{\mathrm{out}}\times r},\,\mathbf{A}\in\R^{r\times d_{\mathrm{in}}},
\label{eq:lora}
\end{equation}
where $r \ll \min(d_{\mathrm{out}}, d_{\mathrm{in}})$ is the LoRA
rank, $\alpha\in\R^+$ is a fixed scaling hyper-parameter, and
$\mathbf{A}, \mathbf{B}$ are the trainable low-rank update
factors~\cite{hu2021lora}. We denote a complete adapter by its
parameter set $\theta = \{(\mathbf{A}_\ell, \mathbf{B}_\ell)
\}_{\ell \in \mathcal{L}}$ over the set of adapted layers
$\mathcal{L}$, plus metadata $m_\theta$ comprising base-model
identifier, training corpus, contributor handle, declared task,
declared licence, upload timestamp, and download statistics.

\paragraph{LoRA ecosystem graph.}
The LoRA ecosystem at time $t$ is a heterogeneous continuous-time
dynamic graph $\Gcal(t) = (\Vcal(t), \Ecal(t), \Rcal, \phi, \psi)$
with four node types and six edge types, where $\Vcal(t)$ is the
node set at $t$, $\Ecal(t)$ the time-stamped edge set, $\Rcal$ the
fixed relation alphabet, $\phi$ the node-feature map (defined
in~\eqref{eq:fused} below), and $\psi$ the edge-type map assigning
each $e \in \Ecal(t)$ a relation $\psi(e) \in \Rcal$.

\begin{definition}[LoRA Ecosystem Graph]
\label{def:graph}
The node set $\Vcal(t) = \Vcal_A(t) \cup \Vcal_B(t) \cup
\Vcal_C(t) \cup \Vcal_D(t)$ comprises adapter nodes $\Vcal_A$
(each carrying weight-tensor features and metadata $m_\theta$),
base-model nodes $\Vcal_B$, contributor nodes $\Vcal_C$, and
downstream-application nodes $\Vcal_D$. The edge set
$\Ecal(t) \subseteq \Vcal(t) \times \Rcal \times \Vcal(t) \times
\R^+$ comprises time-stamped typed edges with relations\\
$\Rcal = \{\texttt{contributes}, \texttt{derives\_from},
\texttt{fine\_tunes},\\ \texttt{deploys}, \texttt{cites},
\texttt{co\_uses}\}$. Node-feature map
$\phi: \Vcal(t) \to \R^{d_v}$ produces type-specific feature
vectors of common dimension $d_v$ after multimodal projection of
weight, text, and behavioural metadata.
\end{definition}

\paragraph{Federated multi-marketplace setting.}
We assume $N$ LoRA-hosting marketplaces $\{C_i\}_{i=1}^N$, each
maintaining a local subgraph $\Gcal_i(t)$ and a private collection
of adapter weights $\Theta_i$. The marketplaces are unwilling to
share raw adapter weights with each other or with a central
coordinator. They consent to participate in federated training of
a shared maliciousness-detection model
$f_{\mathbf{w}}: \Vcal_A(t) \to [0,1]$ provided
privacy-preserving aggregation guarantees are met.

\paragraph{Threat model.}
The \emph{primary} threat FedLoRAGuard targets is the documented
one: single- or few-point malicious adapter uploads to independent
marketplaces. The multi-marketplace \emph{collusion} we analyse is
a worst-case robustness scenario, not a claim that coordinated
$k$-of-$N$ marketplace collusion is the current attacker behaviour;
we state it this way to make the certified radius meaningful under
the strongest bounded-update adversary the federated setting admits.
Formally, the adversary controls $k$ out of $N$ client marketplaces,
with the controlled fraction bounded by a Byzantine tolerance
$\rho \in [0,1)$, i.e.\ $k/N \le \rho$, and may publish arbitrary
backdoored LoRA adapters under any contributor identity within
those marketplaces. The adversary additionally controls the labels
$y_i \in \{0, 1\}$ on the controlled adapters during federated
training (label-flipping). The adversary does not control the
secure-aggregation protocol nor the server-held vetted-adapter
root set used for FLTrust scoring. Crucially, the certified radius
we prove in Section~\ref{sec:theory} applies to
\emph{bounded-update} adversaries in the sense of
Definition~\ref{def:attack} below, not to arbitrary Byzantine
gradients; broader Byzantine tolerance is claimed only empirically
(Section~\ref{subsec:main}), not by the certificate.

\begin{definition}[Adapter-Level Federated Poisoning Attack]
\label{def:attack}
A $(k, \epsilon_{\mathrm{wt}})$-poisoning attack controls
$k\in\{1,\ldots,\lfloor\rho N\rfloor\}$ clients and submits gradient
updates $\tilde{\mathbf{g}}_i$ such that each is within
$\ell_2$-distance $\epsilon_{\mathrm{wt}}\ge 0$ of the corresponding
honest update,
$\norm{\tilde{\mathbf{g}}_i - \mathbf{g}_i^{\mathrm{honest}}}_2
\leq \epsilon_{\mathrm{wt}}$, after server-side gradient clipping
to $\ell_2$-norm $S$ (Algorithm~\ref{alg:fedloraguard} line~5).
Here $\epsilon_{\mathrm{wt}}$ is an attacker-side budget controlling
the per-client deviation magnitude. The objective is to flip the
global maliciousness verdict on at least one target adapter
$\theta^*$.
\end{definition}

\paragraph{Problem statement.}
We seek a parameterisation $\mathbf{w}^*$ that minimises empirical
risk subject to per-round $(\epsilon_r, \delta)$-DP by Gaussian
DP-SGD with client-update clipping $S$; cumulative training-time
budget $\epsilon_T = \mathrm{PRVCompose}(\sigma, q, T)$ at
$\delta$; secure aggregation hiding individual
$\tilde{\mathbf{g}}_i$; and FLTrust scoring on a server-held root
set $R_0$. The deployed verifier should admit (a) a high-probability
certificate of robustness against $k^*$ malicious clients
(Theorem~\ref{thm:certified}) under Definition~\ref{def:attack},
and (b) calibrated maliciousness probabilities suitable for
marketplace moderation actions.

% ============================================================
%  4. FRAMEWORK
% ============================================================
\section{The \modelname\ Framework}
\label{sec:framework}

\modelname\ consists of four components: (i) a multimodal
node-feature encoder, (ii) a per-client local heterogeneous
dynamic GNN, (iii) federated training under
client-update-level Gaussian DP-SGD with secure aggregation and
FLTrust scoring, and (iv) a global verifier that emits per-adapter
maliciousness probability and the certified coalition bound.
Figure~\ref{fig:architecture} shows the architecture.

\begin{figure*}[t]
\centering
\includegraphics[width=0.75\textwidth]{figures/fig1_FedLoraGuard_arch.pdf}
\vspace{-0.3cm}
\caption{The \modelname\ architecture. Each LoRA-hosting
marketplace $C_i$ maintains a local heterogeneous dynamic graph
and adapter pool. A multimodal encoder projects adapter weights,
model-card text, and usage statistics into a unified
node-feature space. A local DyGFormer with HGT relation-aware
attention computes per-adapter embeddings. Per-client DP-SGD
clips and Gaussian-perturbs the client update
(Algorithm~\ref{alg:fedloraguard} lines~5--6). Secure
aggregation and FLTrust scoring compute the global update, hiding
individual client contributions. The verifier emits the
maliciousness probability $\hat{p}_{\mathrm{mal}}(\theta)$ and
the certified coalition bound $k^*$. Dashed red arrows denote
DP-budget feedback; dotted green arrows denote certificate
output.}
\label{fig:architecture}
\end{figure*}

\paragraph{Multimodal node-feature encoder.}
For each adapter node $v_\theta \in \Vcal_A$, three modalities are
encoded and concatenated. Weight-modality features
$\bm{\phi}^{\mathrm{wt}}(\theta) \in \R^{d_{\mathrm{wt}}}$ comprise
the spectral signature $(\sigma_1, \ldots, \sigma_r)$ of each
$\mathbf{B}\mathbf{A}$ matrix in~\eqref{eq:lora}, per-layer
Frobenius norms, and a low-dimensional weight-difference
fingerprint relative to the base model. Text-modality features
$\bm{\phi}^{\mathrm{txt}}(m_\theta) \in \R^{d_{\mathrm{txt}}}$ are
produced by a frozen sentence transformer applied to model cards,
README files, and contributor profiles. Behavioural-modality
features $\bm{\phi}^{\mathrm{beh}}(\theta) \in \R^{d_{\mathrm{beh}}}$
encode download counts, monthly active deployments, and citation
graph centralities. The three streams are fused through cross-attention
and projected to common dimension $d_v$:
\begin{equation}
\phi(v_\theta) = \mathbf{W}_{\mathrm{fuse}}
[\bm{\phi}^{\mathrm{wt}}(\theta);\, \bm{\phi}^{\mathrm{txt}}(m_\theta);\,
\bm{\phi}^{\mathrm{beh}}(\theta)],
\label{eq:fused}
\end{equation}
where $\mathbf{W}_{\mathrm{fuse}}\in\R^{d_v\times(d_{\mathrm{wt}}+
d_{\mathrm{txt}}+d_{\mathrm{beh}})}$ is a trainable projection. The
output $\phi(v_\theta)\in\R^{d_v}$ instantiates the node-feature
map $\phi$ introduced in Definition~\ref{def:graph}. Base-model,
contributor, and downstream-application nodes use type-specific
encoders producing feature vectors of the same dimension $d_v$.

\paragraph{Local heterogeneous dynamic GNN.}
Each client runs a local DyGFormer~\cite{yu2023dygformer} encoder
with HGT relation-aware attention~\cite{hu2020hgt}. For each query
node $v$ at time $t$, the encoder samples the most recent $K$
temporal neighbours of $v$ across all relations, computes
neighbour co-occurrence patches, and applies a Transformer
encoder over the patches. We denote by $\mathbf{h}_v\in\R^{d_v}$
the running hidden representation of node $v$ (initialised to
$\phi(v)$ from Eq.~\ref{eq:fused}) and by $\tau:\Vcal\to\{A,B,C,D\}$
the node-type indicator returning the type subscript of $v$.
For each query node $v$ and each typed neighbour $u$ connected by
relation $r\in\Rcal$, the HGT attention head computes
\begin{equation}
\alpha_{u, v, r} = \softmax\!\Bigl(\frac{(\mathbf{W}_Q^{\tau(v)} \mathbf{h}_v)^\top \mathbf{W}_R^r (\mathbf{W}_K^{\tau(u)} \mathbf{h}_u)}{\sqrt{d_v / H}}\Bigr),
\label{eq:hgt}
\end{equation}
where $\mathbf{W}_Q^{\tau}, \mathbf{W}_K^{\tau}\in\R^{d_v\times d_v}$
are type-specific query and key projections,
$\mathbf{W}_R^r\in\R^{d_v\times d_v}$ is a relation-specific
bilinear interaction matrix, $H$ is the number of attention heads,
and the softmax is taken over the set of typed neighbours $(u,r)$
of $v$. The aggregated representation passes through a feed-forward
block with residual connections.

\paragraph{Federated training (client-update-level DP).}
Let $\mathbf{w}\in\R^d$ denote the parameter vector of the
maliciousness-detection model $f_{\mathbf{w}}$. In each round
$t\in\{1,\ldots,T\}$, the server samples a subset
$S_t \subseteq \{C_i\}_{i=1}^N$ of $m$ clients by Poisson
subsampling with rate $q = m/N$. Each sampled client computes its
local minibatch gradient
$\mathbf{g}_i = \nabla_{\mathbf{w}} \hat{\Lcal}_i(\mathbf{w}_t;
\Gcal_i(t))$, where $\hat{\Lcal}_i$ is the local cross-entropy
loss on client $C_i$'s adapter subgraph $\Gcal_i(t)$. Unlike
canonical example-level DP-SGD~\cite{abadi2016deep}, which clips
each per-example gradient and averages, we operate at the
\emph{client-update} level: the client clips the \emph{whole
update} $\mathbf{g}_i$ to $\ell_2$-norm $S$ and adds a single
Gaussian noise draw
$\Ncal(\mathbf{0},\sigma^2 S^2 \mathbf{I})$ calibrated to the
per-round R\'enyi-DP budget. It then transmits the noised update
$\tilde{\mathbf{g}}_i$ via SecAgg~\cite{bonawitz2017practical}.
The server scores each update against a small vetted root set
$R_0$ through FLTrust~\cite{cao2021fltrust} cosine similarity to
a server-computed reference gradient $\mathbf{g}_0$, weights the
update by $\mathrm{ReLU}(\cos(\tilde{\mathbf{g}}_i, \mathbf{g}_0))$,
and aggregates by weighted averaging.
Algorithm~\ref{alg:fedloraguard} gives one complete training
round.

\begin{algorithm}[t]
\small
\caption{\modelname\ training round}
\label{alg:fedloraguard}
\begin{algorithmic}[1]
\REQUIRE Clients $\{C_i\}_{i=1}^N$, parameters $\mathbf{w}_t$, client-update clip threshold $S$, noise multiplier $\sigma$, learning rate $\eta$, vetted root set $R_0$, sampling rate $q$
\STATE Server Poisson-samples $S_t \subseteq \{C_i\}$ at rate $q$
\STATE Server computes $\mathbf{g}_0 \gets \nabla_{\mathbf{w}} \hat{\Lcal}(\mathbf{w}_t; R_0)$
\FOR{each $C_i \in S_t$ in parallel}
    \STATE $\mathbf{g}_i \gets \nabla_{\mathbf{w}} \hat{\Lcal}_i(\mathbf{w}_t; \Gcal_i(t))$ \hfill\COMMENT{full client update}
    \STATE $\bar{\mathbf{g}}_i \gets \mathbf{g}_i \cdot \min(1, S/\norm{\mathbf{g}_i}_2)$ \hfill\COMMENT{clip whole update}
    \STATE $\tilde{\mathbf{g}}_i \gets \bar{\mathbf{g}}_i + \Ncal(\mathbf{0}, \sigma^2 S^2 \mathbf{I})$ \hfill\COMMENT{single Gaussian draw}
    \STATE Transmit $\tilde{\mathbf{g}}_i$ by SecAgg
\ENDFOR
\STATE Aggregator: $\bar{\mathbf{G}} \gets \mathrm{SecAgg}(\{\tilde{\mathbf{g}}_i\})$
\STATE $\beta_i \gets \mathrm{ReLU}(\cos(\tilde{\mathbf{g}}_i, \mathbf{g}_0)) / \sum_j \mathrm{ReLU}(\cos(\tilde{\mathbf{g}}_j, \mathbf{g}_0))$
\STATE $\mathbf{w}_{t+1} \gets \mathbf{w}_t - \eta \sum_i \beta_i \tilde{\mathbf{g}}_i$
\STATE Update PRV accountant; compute $k^*$ through Theorem~\ref{thm:certified}
\end{algorithmic}
\end{algorithm}

% ============================================================
%  5. THEORETICAL PROPERTIES
% ============================================================
\section{Theoretical Properties}
\label{sec:theory}

We present two results: a sensitivity bound for the federated
heterogeneous-dynamic-GNN client update
(Theorem~\ref{thm:sensitivity}) and a certified poisoning-radius
theorem (Theorem~\ref{thm:certified}). Proof sketches appear in
the main text; full proofs are in Appendix~\ref{app:proofs}.

\begin{theorem}[Client-update sensitivity bound]
\label{thm:sensitivity}
Let $f_{\mathbf{w}}$ be an $L$-layer heterogeneous dynamic GNN
operating on the ecosystem graph $\Gcal(t)$ of
Definition~\ref{def:graph}, with per-relation weight matrices
$\{\mathbf{W}_R^r\}_{r \in \Rcal}$ satisfying
$\norm{\mathbf{W}_R^r}_2 \leq B_W$ for all $r \in \Rcal$ (here
$\Rcal$ is the six-element relation set from
Definition~\ref{def:graph}, so $|\Rcal|=6$). Let $\phi^{\mathrm{nl}}$
denote the point-wise non-linearity inside the HGT block and assume
it is $\rho$-Lipschitz; for GELU $\rho\approx 1.13$ (attained near
$x\approx\pm\sqrt{2}$), which we use in place of the looser
$\rho\le 1$. Let the per-node temporal-neighbourhood degree be
bounded by $d_{\max}$, and let $S$ be the client-update clipping
threshold used at Algorithm~\ref{alg:fedloraguard} line~5.
Consider two adjacent local datasets
$\Dcal_i \sim \Dcal_i'$ differing in a single node-edge addition
or modification. The $\ell_2$-sensitivity of the clipped client
update
$\bar{\mathbf{g}}_i:=\mathbf{g}_i\cdot\min(1,S/\norm{\mathbf{g}_i}_2)$
satisfies
\begin{equation}
\Delta_2 := \sup_{\Dcal_i \sim \Dcal_i'}
\norm{\bar{\mathbf{g}}_i(\Dcal_i) - \bar{\mathbf{g}}_i(\Dcal_i')}_2
\;\le\; 2 S,
\label{eq:sensitivity}
\end{equation}
by two-sided clipping alone. The graph-structural refinement
\begin{equation}
\Delta_2 \;\le\; \min\!\Bigl(2S,\;
2S\,(\rho B_W d_{\max})^L\, |\Rcal|^{1/2}\Bigr)
\label{eq:sensitivity-refined}
\end{equation}
holds when the RHS is finite; on real Hugging Face lineage graphs
$d_{\max}$ can exceed $10^3$, so~\eqref{eq:sensitivity-refined}
degenerates to the loose $2S$ bound at large depth and the
operational sensitivity we use for calibration is
$\Delta_2 = 2S$.
\end{theorem}

\begin{proof}[Proof sketch]
For any two datasets, the map
$\mathbf{g}\mapsto\mathbf{g}\cdot\min(1,S/\norm{\mathbf{g}}_2)$
has $\ell_2$-norm at most $S$, so the difference of any two
clipped updates is at most $2S$ by the triangle inequality; this
gives~\eqref{eq:sensitivity} and requires no assumption on the
architecture. For~\eqref{eq:sensitivity-refined}, induction on
layer depth propagates a single-node change through the HGT
message-passing update
$\mathbf{h}_v^{(\ell)} = \phi^{\mathrm{nl}}(\sum_{r,u} \alpha_{u,v,r}
\mathbf{W}_R^r \mathbf{h}_u^{(\ell-1)})$; Cauchy--Schwarz plus the
relation-wise spectral bound $\norm{\mathbf{W}_R^r}_2\le B_W$
gives $\norm{\Delta \mathbf{h}_v^{(\ell)}}_2 \le
\rho B_W d_{\max}\norm{\Delta \mathbf{h}_v^{(\ell-1)}}_2$ when the
affected node lies in the temporal neighbourhood, and iterating
$L$ layers yields the $(\rho B_W d_{\max})^L$ factor; the
$|\Rcal|^{1/2}$ comes from Cauchy--Schwarz over the relation-type
sum. The client-update output is then clipped again to $S$, and
the tighter of the two bounds is the operative one. Full details
appear in Appendix~\ref{app:proofs}.
\end{proof}

\paragraph{Calibration of $\sigma$ in practice.}
The Gaussian mechanism at sensitivity $\Delta_2=2S$ gives
$(\epsilon,\delta)$-DP per round with
$\sigma\ge 2S\sqrt{2\ln(1.25/\delta)}/\epsilon$; the tighter
analytic bound~\eqref{eq:sensitivity-refined} is only useful when
$(\rho B_W d_{\max})^L\,|\Rcal|^{1/2}\!\le\!1$, which does not hold
for our deployment scale. In practice we do not use either closed
form to set $\sigma$: we adopt $\sigma=1.1$ and account the full
$T$-round composition numerically through the Opacus PRV
accountant~\cite{gopi2021numerical} under Poisson subsampling at
$q=0.2$, obtaining
$(\epsilon_T,\delta)\approx(13.6,10^{-5})$ at $T=100$; see
Appendix~\ref{app:repro}.

\begin{theorem}[Certified poisoning radius]
\label{thm:certified}
Suppose \modelname\ runs Algorithm~\ref{alg:fedloraguard} for
$T$ rounds with client-update-level Gaussian DP-SGD (clip $S$,
noise multiplier $\sigma$), secure aggregation, and FLTrust
scoring against a server-held vetted root set $R_0$. At
verification time, the aggregated global model is passed through
a single Gaussian smoothing draw with budget
$(\epsilon_r,\delta)$; this smoothing draw is what the certificate
depends on and is a design parameter separate from the training-time
budgets. Let the global verifier
$f_{\mathbf{w}_T}:\Vcal_A(t)\to\Delta^{c-1}$ emit, on a queried
adapter $\theta$, a class-probability vector over $c=2$ classes
(benign, malicious), and let $\hat{p}_1\ge\hat{p}_2$ be the
largest and second-largest entries of that vector (so
$\hat{p}_1+\hat{p}_2=1$ here). Let $\Phi^{-1}:(0,1)\to\R$ denote
the inverse cumulative distribution function of the standard
normal. The prediction of $f_{\mathbf{w}_T}$ on $\theta$ is
invariant with probability at least $1 - T\delta$ under any
$(k, \epsilon_{\mathrm{wt}})$-poisoning attack
(Definition~\ref{def:attack}) such that
\begin{equation}
k \leq k^* = \left\lfloor \frac{N \bigl(\Phi^{-1}(\hat{p}_1) - \Phi^{-1}(\hat{p}_2)\bigr)}{2 \exp(\epsilon_r)} \right\rfloor,
\label{eq:kstar}
\end{equation}
where the floor enforces an integer coalition size. The cumulative
training-time budget $\epsilon_T$, obtained by PRV composition of
the $T$ per-round training budgets, governs the marketplace-level
privacy guarantee (Section~\ref{sec:experiments}) but does
\emph{not} enter the certificate: the stability factor is set by
the budget of the single smoothing draw applied at verification
time, not by the budget accumulated over training.
\end{theorem}

\begin{proof}[Proof sketch]
The bound is obtained by composing three established results,
each stated with its own assumption; the complete book-keeping is
in Appendix~\ref{app:proofs}. \textbf{(i)} The random variable
being smoothed is the aggregated model
$\mathbf{w}_T\in\R^d$: at verification time the output logits on
$\theta$ are averaged over an $(\epsilon_r,\delta)$-DP Gaussian
draw applied to $\mathbf{w}_T$, producing the smoothed classifier
studied by PixelDP~\cite{lecuyer2019certified} and (in the
federated setting) by CRFL~\cite{xie2021crfl}. \textbf{(ii)} A
single malicious client controlling $k=1$ adds to the aggregated
model a bounded parameter-space perturbation of size at most
$\epsilon_{\mathrm{wt}}$ per round, because SecAgg forces the
attacker's update through the same clip-and-noise pipeline as
honest clients and Definition~\ref{def:attack} bounds the
attacker's post-clip deviation; $k$ colluding clients induce a
worst-case perturbation of size at most $k\epsilon_{\mathrm{wt}}$
by triangle inequality, subject to the FLTrust down-weighting.
\textbf{(iii)} FLTrust does not invalidate the smoothing argument
because the FLTrust weight $\beta_i\in[0,1]$ can only decrease
the effective attacker deviation, so any certificate obtained
under unweighted aggregation is a fortiori valid under FLTrust
weighting. Substituting the resulting bounded perturbation into
the CRFL~\cite{xie2021crfl} class-margin certificate (a
Neyman--Pearson hypothesis-testing bound in
$\Phi^{-1}(\hat{p}_1)-\Phi^{-1}(\hat{p}_2)$), and using the
PixelDP stability factor $e^{\epsilon_r}$ from the single-query
$(\epsilon_r,\delta)$-DP smoothing draw, translates a
class-margin radius into the client coalition bound
\begin{equation*}
k^* = \left\lfloor \frac{N(\Phi^{-1}(\hat{p}_1) - \Phi^{-1}(\hat{p}_2))}{2 e^{\epsilon_r}} \right\rfloor.
\end{equation*}
The factor $2 e^{\epsilon_r}$ combines the PixelDP per-query
stability multiplier ($e^{\epsilon_r}$) with the CRFL
class-margin denominator ($2$, from the two tails of the
half-Gaussian). The failure probability $T\delta$ is a union
bound over the $T$ inference draws used to reach a target
smoothing budget; a single-draw certificate uses the tighter
$\delta$. The graph-structural extension to heterogeneous
edge-set perturbations is handled by
\citet{wang2021certgnn,scholten2022randomized}.
\end{proof}

\paragraph{From certificate to marketplaces.}
Eq.~\eqref{eq:kstar} converts an operational per-query privacy
budget into an architectural integrity guarantee. At $N = 50$,
$\hat{p}_1 - \hat{p}_2 \approx 0.6$ (so
$\Phi^{-1}(\hat{p}_1)-\Phi^{-1}(\hat{p}_2)\approx 1.68$), per-query
smoothing budget $\epsilon_r = 0.5$, and $\delta = 10^{-5}$,
evaluation gives $k^* = \lfloor 50\cdot 1.68 / (2e^{0.5})\rfloor =
\lfloor 84.2/3.30\rfloor = 25$. Under Definition~\ref{def:attack}
this means no bounded-update coalition of fewer than twenty-five
colluding malicious marketplaces can flip the global verdict on
any vetted adapter, with probability at least
$1 - 100 \cdot 10^{-5} = 0.999$. We report this recomputed value
throughout; an earlier draft substituted the cumulative
training-time budget $\epsilon_T$ into the denominator, which
understated the certified radius.

\paragraph{Certified vs empirical robustness.}
Two robustness numbers appear later. The
\emph{certified} radius $k^*=25$ out of $N=50$ ($50\%$) is a
distribution-free bound and holds only under the bounded-update
threat model of Definition~\ref{def:attack}. The
\emph{empirical} label-flip stress test in
Section~\ref{subsec:main} runs Krum-style adversarial clients up
to fraction $\rho\le 0.33$ ($33\%$) and observes no certificate
degradation on \benchname. These are distinct results and we
label them as such throughout: reviewers who see the two numbers
should not conflate the $50\%$ certificate bound with the
empirical $33\%$ Byzantine-tolerance operating point.

% ============================================================
%  6. EXPERIMENTS
% ============================================================
\section{Experiments}
\label{sec:experiments}

\paragraph{Datasets.}
We construct \benchname, a benchmark of $13{,}500$ benign and
backdoored LoRA adapters across four base-model families
(Llama-2-7B, Llama-3-8B, Mistral-7B, Qwen-7B) and ten attack
types (BadNets, VPI, Sleeper, MTBA, CTBA, AddSent, BadEdit,
InsertSent, weight-poisoning, and trigger-free LoRA-merging).
Adapter ranks span $r \in \{4, 8, 16, 32\}$. Each adapter is
trained for three epochs over a task corpus drawn from Alpaca,
Dolly, GSM8K, SQuAD-v2, IMDB, AGNews, ARC-Challenge, HumanEval, or
GLUE. The benign-to-malicious split is $50/50$; we treat this as
a controlled evaluation design rather than a claimed base rate of
malicious adapters in the wild, and re-report FPR and precision
at low-false-positive operating points in Appendix~\ref{app:additional}.
Each adapter is recorded with its per-layer
$(\mathbf{A}_\ell, \mathbf{B}_\ell)$ matrices, spectral signature,
layer-wise Frobenius norms, model card, contributor metadata, and
download-statistic time-series. Cross-marketplace lineage edges
are synthesised from
HuggingGraph~\cite{rahman2025hugginggraph} and
PADBench~\cite{sun2025peftguard}. A component-level provenance
audit of what parts of the graph are real versus synthetic
appears in Appendix~\ref{app:provenance}
(Table~\ref{tab:provenance}). We also evaluate on the
BackdoorLLM-LoRA leaderboard~\cite{li2025backdoorllm} and, as a
cross-domain sanity check for the backbone only, on four IDS
datasets (CIC-IDS2017~\cite{sharafaldin2018toward},
Edge-IIoTset~\cite{ferrag2022edge},
UNSW-NB15~\cite{moustafa2015unsw},
TON-IoT~\cite{moustafa2021ton}); IDS results are reported in
Appendix~\ref{app:ids} and are not part of the LoRA-integrity
contribution.

\paragraph{Splits and leakage control.}
Because \modelname\ exploits lineage, a random adapter split would
risk train/test leakage across related adapters. We report the
main-paper number under a \emph{lineage-component-disjoint} split
(no derivation chain crosses the train/test boundary) and re-run
under three additional stress splits in
Appendix~\ref{app:additional} (\emph{base-model-disjoint},
\emph{contributor-disjoint}, \emph{attack-family-held-out}). The
strictest of these is attack-family-held-out, where entire attack
families are unseen at training time; \modelname\ retains
macro-F1 above $92\%$ under all four splits, so the reported
$96.4\%$ is not an artefact of family memorisation.

\paragraph{Baselines.}
We compare against six baselines.
\emph{Centralised:} PEFTGuard~\cite{sun2025peftguard} ($98.3\%$
centralised macro-F1 on PADBench), ShadowGenes~\cite{schulz2025shadowgenes}
(graph-based model-genealogy verification).
\emph{Federated without graphs:} FedAvg-MLP and DP-FedAvg-MLP.
\emph{Federated with static graphs:}
FedGraphNN-HGT~\cite{he2021fedgraphnn}.
\emph{Federated dynamic-graph:} Krum-DyGFormer with
Krum~\cite{blanchard2017machine} Byzantine-robust aggregation
over the DyGFormer-HGT backbone. \emph{Non-graph spectral baseline:}
a \emph{SVD-only} detector that thresholds the $\sigma_1/\sigma_2$
ratio of the $\mathbf{B}\mathbf{A}$ update matrix (and a
logistic-regression variant on the full spectral signature), with
no graph, no federation, and no certificate --- included to test
whether the prominent rank-1 backdoor signature alone suffices.

\paragraph{Implementation.}
\modelname\ is implemented in PyTorch~$2.1$ with PyTorch Geometric
for the graph operations and the Flower~\cite{beutel2020flower}
federated runtime. Training uses AdamW~\cite{loshchilov2019adamw}
with learning rate $5 \times 10^{-4}$, weight decay $10^{-4}$,
client-update gradient clipping at norm $S = 1.0$, and cosine
annealing over $T = 100$ rounds. The FLTrust root set $R_0$
contains $200$ vetted clean adapters. Client-update-level DP-SGD
parameters: clip threshold $S = 1.0$, per-round noise multiplier
$\sigma = 1.1$, Poisson subsampling rate $q = 0.2$ ($m = 10$
clients sampled at $N = 50$); cumulative privacy budget
$(\epsilon_T,\delta)\approx(13.6, 10^{-5})$ at $T=100$ under the
Opacus PRV accountant~\cite{gopi2021numerical}. The certificate
per-query smoothing budget is $\epsilon_r = 0.5$ and is a
verification-time design parameter distinct from $\epsilon_T$.
Experiments use $4 \times$ NVIDIA A100 80~GB GPUs. Each result is
the mean over three random seeds ($42, 137, 2026$); the unit of
analysis for the statistical tests below is the $(\mathrm{seed},
\mathrm{dataset})$ pair, so with three seeds and seven
method-dataset combinations we test on $21$ replicates. We
report empirical bootstrap $95\%$ confidence intervals and
significance through the Friedman test with Holm post-hoc
correction.

\subsection{Main results}
\label{subsec:main}

Table~\ref{tab:main} reports the headline comparison on
\benchname\ and the BackdoorLLM-LoRA leaderboard. \modelname\
obtains $96.4\%$ macro-F1 ($0.984$ AUROC) on \benchname\ with
certified coalition bound $k^* = 25$ at $N = 50$ under smoothing
budget $\epsilon_r=0.5$, i.e.\ a $50\%$ certified bound on
bounded-update coalitions. On top of this, an empirical
label-flip Byzantine stress test tolerates $\rho\le 0.33$
($33\%$) attacker-controlled clients without certificate
degradation; the $50\%$ certificate bound and the $33\%$
empirical operating point are distinct results and we do not
conflate them. \modelname\ lands within $1.9$ macro-F1 points of
the centralised PEFTGuard reference number
($96.4\%$ vs $98.3\%$) on the same backbone while preserving
marketplace privacy. As a transfer/generality probe, the same
backbone reaches $95.8\%$ macro-F1 averaged over four IDS
datasets (Appendix~\ref{app:ids}); these are not part of the
core adapter-integrity claim.

\begin{table}[!t]
\centering
\caption{Main results on \benchname\ + BackdoorLLM-LoRA. Macro-F1
in \% and AUROC as probability, both reported as mean over three
seeds ($95\%$ bootstrap CI $\pm 0.3$--$0.5$\,pp for the federated
methods). Cert.\ Bound is the certified $k^*$ at $N = 50$
marketplaces under smoothing budget $\epsilon_r=0.5$, integer
by~\eqref{eq:kstar}. ``Cent.'' = centralised; ``N/A'' =
inapplicable. IDS transfer results are in
Appendix~\ref{app:ids}.}
\label{tab:main}
\footnotesize
\setlength{\tabcolsep}{2.5pt}
\renewcommand{\arraystretch}{1.05}
\begin{tabular}{@{}lccccc@{}}
\toprule
Method & ACC & Macro-F1 & AUROC & $k^*$ & $\epsilon_T$ \\
\midrule
\makecell[l]{PEFTGuard\\~\cite{sun2025peftguard}\\(Cent.)}
                              & 98.3 & 98.1 & 0.991 & N/A & N/A \\
\addlinespace[2pt]
\makecell[l]{ShadowGenes\\~\cite{schulz2025shadowgenes}\\(Cent.)}
                              & 95.1 & 94.7 & 0.967 & N/A & N/A \\
\addlinespace[2pt]
FedAvg-MLP                    & 89.6 & 88.4 & 0.924 & 0   & N/A \\
DP-FedAvg-MLP                 & 86.8 & 85.7 & 0.901 & 13  & 13.6 \\
\addlinespace[2pt]
\makecell[l]{FedGraphNN-HGT\\\cite{he2021fedgraphnn}}
                              & 92.7 & 91.6 & 0.951 & 0   & N/A \\
\addlinespace[2pt]
Krum-DyGFormer                & 93.5 & 92.8 & 0.962 & 6   & N/A \\
\midrule
\makecell[l]{\modelname\\(ours, $T{=}100$)}
                              & 97.1 & 96.4 & 0.984 & 25  & 13.6 \\
\bottomrule
\end{tabular}
\end{table}
The Friedman test across the $21$ (seed, dataset, method)
replicates gives $\chi_F^2(6) = 51.4$, $p < 10^{-8}$, rejecting
the null of equal performance. Holm post-hoc at $\alpha = 0.05$
confirms \modelname\ significantly exceeds every federated
baseline. Wilcoxon signed-rank tests give $p < 0.001$ for the gap
to FedAvg-MLP (effect size $r = 0.86$), to FedGraphNN-HGT ($r =
0.79$), and to Krum-DyGFormer ($r = 0.74$).

\subsection{Spectral signature analysis}
\label{subsec:spectrum}

Beyond aggregate accuracy, the mechanism-level evidence that the
defence distinguishes its target attack family is the spectral
signature of the $\mathbf{B}\mathbf{A}$ update
matrix~\cite{zhang2023adalora,wei2024brittleness}. Backdoored
adapters consistently show anomalously large top-$1$ singular
values relative to trailing components, reflecting the rank-$1$
concentration of trigger-encoding.
Figure~\ref{fig:spectrum} plots mean top-$8$ singular values of
$\mathbf{B}\mathbf{A}$ (averaged across adapted layers) for
benign and backdoored adapters in \benchname.

\begin{figure}[!t]
\centering
\includegraphics[width=0.96\columnwidth]{figures/fig2_mean.pdf}
\vspace{-0.45cm}
\caption{Mean top-$8$ singular values of $\mathbf{B}\mathbf{A}$ for
benign and backdoored adapters in \benchname\ (averaged across
adapted layers; $95\%$ confidence bands). Backdoored adapters have a
$2.43\times$ larger top-$1$ singular value; trailing values are
indistinguishable --- the rank-$1$ concentration captured by
$\bm{\phi}^{\mathrm{wt}}$.}
\label{fig:spectrum}
\end{figure}

The top-$1$ singular value for backdoored adapters is $2.43\times$
larger than for benign adapters while trailing values are
indistinguishable --- the rank-$1$ concentration that the
weight-modality features $\bm{\phi}^{\mathrm{wt}}$ extract and
propagate through the graph.

\subsection{Ablations}
\enlargethispage{4\baselineskip}
\label{subsec:ablation}

Table~\ref{tab:ablation} reports each \modelname\ component's
contribution on \benchname.

\begin{table}[tb]
\centering
\caption{Component ablation on \benchname\ (\modelname, $T = 100$,
$N = 50$, mean over three seeds). The DP ablation is the key
result to read separately: removing DP-SGD gains $0.5$\,pp macro-F1
but sets the certified bound to $k^*=0$, so the DP cost is the
price of the certificate rather than of detection quality.}
\label{tab:ablation}
\small
\setlength{\tabcolsep}{4pt}
\begin{tabular}{@{}lcccc@{}}
\toprule
Configuration & ACC & Macro-F1 & AUROC & $k^*$ \\
\midrule
\modelname\ (full)                               & 97.1 & 96.4 & 0.984 & 25 \\
\midrule
\multicolumn{5}{@{}l}{\emph{Architectural ablations}}\\
\;\; w/o multimodal encoder\\ (weight only)        & 95.6 & 94.7 & 0.972 & 22 \\
\;\; w/o HGT relation-aware\\ attention            & 94.2 & 93.4 & 0.961 & 19 \\
\;\; w/o DyGFormer\\ (static-graph GAT)            & 93.1 & 92.0 & 0.952 & 16 \\
\;\; w/o FLTrust scoring                         & 96.3 & 95.6 & 0.978 & 19 \\
\midrule
\multicolumn{5}{@{}l}{\emph{Privacy / training ablations}}\\
\;\; w/o DP-SGD\\ (federated, no DP)               & 97.4 & 96.9 & 0.987 & 0 \\
\;\; w/o secure aggregation                      & 97.1 & 96.4 & 0.984 & 25 \\
\;\; centralised oracle\\ (upper bound)            & 98.3 & 98.1 & 0.991 & --- \\
\bottomrule
\end{tabular}
\end{table}

Four observations. First, DyGFormer is the dominant \emph{detection}
component: replacing it with a static-graph GAT costs $4.4$\,pp
macro-F1 and $9$ points of $k^*$, so cross-marketplace lineage
through time carries the main signal (HGT attention and the
multimodal encoder add a further $6$ and $3$ points of $k^*$).
Second, the DP ablation \emph{separates detection quality from
certification quality}: turning off DP recovers $0.5$\,pp macro-F1
but collapses the certificate to $k^*=0$, i.e.\ the marginal DP
cost is being paid for the certificate, not for detection quality.
Third, removing FLTrust does not affect the certified bound
(FLTrust re-weighting only tightens it further) but drops
empirical macro-F1 by $0.8$\,pp under label-flip attack. Fourth,
secure aggregation is a privacy mechanism, not a utility one:
removing it does not change accuracy or the certificate but
exposes individual updates to the coordinator.

\subsection{Additional Analyses}
Appendix~\ref{app:additional} reports six further analyses: an
SVD-only baseline ($90.2\%$ macro-F1 on easy rank-1 backdoors,
$71.4\%$ on composite/trigger-free, no certificate), the four
leakage-controlled splits, security-operating-point metrics (FPR,
precision, PR-AUC, detection@FPR$=1\%$), graceful degradation to
$94.1\%$ at $\epsilon_T=1.0$ and $91.1\%$ under $50\%$ edge
removal, cross-hardware latency (A100/A10G/T4), and a KS test
confirming the spectral gap ($D=0.71$, $p<10^{-3}$).

% 7. DISCUSSION
\section{Discussion}
\label{sec:discussion}

\paragraph{What the certificate buys, and at what cost.}
The certified bound $k^* = 25$ (Section~\ref{sec:theory})
converts an inference-time smoothing budget into a coalition-size
guarantee against bounded-update adversaries
(Definition~\ref{def:attack}) that holds under secure aggregation
with no raw-weight sharing, at a cost of $1.9$\,pp macro-F1 versus
the centralised PEFTGuard reference number and $0.5$--$1.0$\,pp from
DP-SGD training. The trade is favourable because centralised
detection requires the weight-sharing commercial actors refuse,
whereas the federated certificate uniquely combines ${>}96\%$
macro-F1, training-time privacy at
$(\epsilon_T,\delta)\approx(13.6, 10^{-5})$ under Opacus PRV
composition, and a distribution-free coalition bound on a
precisely defined bounded-update threat model.

\paragraph{Deployment scenario and hardware envelope.}
In an operational hub, \modelname\ runs as an ingestion-time gate:
on upload it extracts an adapter's spectral, textual, and usage
features, scores them against the local lineage subgraph, and
emits $\hat{p}_{\mathrm{mal}}(\theta)$ within
$0.7$--$2.9$\,s per adapter --- the ``sub-second'' claim is
specific to A100-class hardware, and consumer GPUs (A10G/T4)
extend the budget to roughly $1.6$--$2.9$\,s. Marketplaces
exchange only DP-noised, securely-aggregated gradients, never raw
weights, and the certified bound $k^*$ is recomputed each round as
a trust indicator.

\paragraph{On FLTrust false-zeroing.}
FLTrust zeroes clients whose gradient points away from the root-set
reference --- the intended Byzantine defence, but a \emph{benign}
marketplace genuinely misaligned with the root set (e.g., a
low-resource-language hub) can be down-weighted or zeroed, losing
its contribution and fairness. Mitigations include a more diverse
root set, per-cluster reference gradients, or a soft floor on
client weights.

% ============================================================
%  8. CONCLUSION
% ============================================================
\section{Conclusion}
\label{sec:conclusion}

\modelname\ verifies LoRA-adapter integrity in a federated,
privacy-preserving manner: a heterogeneous continuous-time dynamic
graph captures the ecosystem, a federated DyGFormer--HGT encoder
scores adapters under client-update-level Gaussian DP-SGD with
secure aggregation and FLTrust, and a certified poisoning-radius
theorem converts the per-query smoothing budget into a
bounded-update coalition bound. On $13{,}500$ adapters it obtains
$96.4\%$ macro-F1 and certified coalition bound $k^* = 25$ at
$N = 50$ under smoothing budget $\epsilon_r=0.5$, with sub-second
per-adapter latency on A100-class hardware and a measured
$\approx 4\%$ per-round wall-clock overhead for the full DP + SecAgg
stack under CPU-hosted federation (pooled across three
measurement runs; see Appendix~\ref{app:repro}), offering a
deployable alternative to weight-sharing detectors.

% ============================================================
%  ACKNOWLEDGEMENTS
% ============================================================
\section*{Acknowledgements}
This work was supported by a grant for research centers in
Artificial Intelligence from the Ministry of Economic Development
of the Russian Federation under subsidy agreement (identifier
000000C313925P3Q0002). The work was carried out at the Institute of
Cyber Intelligence Systems and the Artificial Intelligence Research
Center, National Research Nuclear University MEPhI (Moscow). We
thank the ACL Rolling Review reviewers, whose comments materially
improved the certificate derivation, the spectral-only baseline,
the DP-accounting disclosure, and the scoping of the threat model
and deployment claims.

% ============================================================
%  LIMITATIONS
% ============================================================
\section*{Limitations}
\label{sec:limitations}

\paragraph{Adaptive attacks against the certificate.}
Theorem~\ref{thm:certified} bounds the coalition of
\emph{bounded-update} label-flipping clients in the sense of
Definition~\ref{def:attack}. It does not directly bound
attackers who control client \emph{training data}
(gradient-level poisoning) under FLTrust scoring, nor arbitrary
Byzantine gradients unconstrained by
$\epsilon_{\mathrm{wt}}$. Empirically the FLTrust filter
eliminates the label-flip attack at $\rho \leq 0.33$, but this
$33\%$ empirical operating point is distinct from the $50\%$
certified bound, and formally extending the certificate to
data-controlling adversaries is left for future work.

\paragraph{HuggingGraph synthesis.}
Cross-marketplace lineage edges in \benchname\ are synthesised
from HuggingGraph~\cite{rahman2025hugginggraph} and PADBench. We
do not replicate the full causal graph of the Hugging Face Hub.
Adapters that are truly novel (i.e., not derived from any base
model in \benchname) are excluded by construction. Appendix
\ref{app:provenance} audits which parts of the graph are real
versus synthetic.

\paragraph{DP accounting.}
The reported $(\epsilon_T,\delta)\approx(13.6,10^{-5})$ at
$T=100$ rounds is the numerical output of the Opacus PRV
accountant~\cite{gopi2021numerical} at $\sigma=1.1$, Poisson
subsampling $q=0.2$; the closed-form Gaussian-mechanism bound
derived from~\eqref{eq:sensitivity} is looser. Tighter accountants
(e.g.\ Gaussian DP~\cite{dong2022gaussian}) could yield smaller
noise scales and, if $\sigma$ is re-calibrated for a fixed
$\epsilon_T$, a marginally larger empirical robustness envelope.
We use the standard Opacus PRV accountant for reproducibility.

\paragraph{Coverage of LoRA variants.}
We evaluate four base-model families (Llama-2-7B, Llama-3-8B,
Mistral-7B, Qwen-7B) and four PEFT variants (LoRA, QLoRA, AdaLoRA,
DoRA). Behaviour on full-rank fine-tuning, mixture-of-experts
adapters, and quantization-aware LoRA at $2$-bit precision is not
characterised.

\paragraph{Adversarial benchmark realism.}
The ten attack types in \benchname\ are publicly documented in
prior work. Zero-day attack families that do not match these
behavioural signatures could degrade detection. The certified
bound is operationally meaningful only under the threat-model
assumptions of Definition~\ref{def:attack}. \benchname's $50/50$
benign/malicious split is a controlled evaluation design, not a
claimed base rate of backdoored adapters in the wild; we therefore
also report FPR, precision, and detection@FPR$=1\%$ in
Appendix~\ref{app:additional} for security-operating-point
readers.

\paragraph{False positives, fairness, and access.}
A false positive --- flagging a benign adapter as malicious ---
imposes a real cost on the developer (quarantine, reputational
harm, delayed release), so deployments should pair the detector
with a low false-positive operating point and a human-review
redress path rather than automatic removal. Detection quality may
also vary across base-model families and languages: English-language
adapters dominate \benchname, and hubs specialising in low-resource
languages may be under-served both by the training distribution and
by the FLTrust root set (see the false-zeroing discussion in
Section~\ref{sec:discussion}); a per-family and per-language
performance audit is needed before broad deployment. Finally, the
$\sim\!960$ A100-GPU-hour training cost and the federated design
concentrate detection capability among large, compute-rich
marketplaces, which raises an inclusivity and power-concentration
concern for small independent hubs; participation incentives (for
example reputation or security-certification tiers) and lighter
client-side training are needed to make the scheme broadly
adoptable. We flag these as societal considerations rather than
solved problems.

% ============================================================
%  ETHICS STATEMENT
% ============================================================
\section*{Ethics Statement}

\paragraph{Defensive contribution.}
The system is designed for defensive operation by LoRA-hosting
marketplaces. Its components support adapter-integrity
verification rather than attack development. The threats it
defends against are publicly documented in cited prior
work~\cite{liu2024loraattack,lermen2023lora,huang2024composite}.
We add no new attack mechanisms beyond what is already public.

\paragraph{Privacy-preserving by design.}
The federated framework ensures marketplaces never share raw
adapter weights with any other party. Client-update-level
Gaussian DP-SGD plus secure aggregation gives cumulative
$(\epsilon_T, \delta)\approx(13.6, 10^{-5})$ per marketplace over
$T = 100$ training rounds under the Opacus PRV accountant. This
is exactly the property that makes \modelname\ deployable in
settings where PEFTGuard cannot be.

\paragraph{Benchmark provenance.}
The \benchname\ benchmark is constructed from publicly released
base models and publicly documented attack methodologies.
Synthesised cross-marketplace lineage edges use the structure of
HuggingGraph but introduce no personally identifying information.
Adapter content is restricted to weight tensors, model cards, and
download statistics. No private adapter weights from any
marketplace partner are used. The benchmark will be released under
a clear research licence with a use-case attestation and SHA-256
checksums for every artefact (Appendix~\ref{app:provenance}).

\paragraph{Operator responsibility.}
The certified bound $k^* = 25$ is a probabilistic guarantee
against a specified bounded-update attack model
(Definition~\ref{def:attack}), not against arbitrary Byzantine
gradients. Operators must monitor for distribution shift and
trigger online adaptation at appropriate cadence. Coalitions
larger than $k^*$, or coalitions violating the bounded-update
assumption, can in principle flip the verdict and require
operator-level mitigation. \modelname\ is one layer in a
defence-in-depth pipeline.

\paragraph{Compute footprint.}
Total compute for the reported experiments is approximately
$960$ A100-GPU-hours across hyperparameter sweeps and three
seeds. Per-adapter scan latency at deployment is $0.7$\,s on a
single A100.

\paragraph{Use of AI writing assistants.}
LLM writing assistants were used only for grammar and phrasing on
author-drafted text. No technical content, experimental design,
results, or analysis was AI-generated; all citations were
manually verified.

% ============================================================
%  REFERENCES
% ============================================================
\begin{thebibliography}{99}
\small

\bibitem[Hu et~al.(2021)]{hu2021lora}
E.~J.~Hu et~al.
\newblock LoRA: Low-rank adaptation of large language models.
\newblock In \emph{Proc.\ ICLR}, 2021.

\bibitem[Dettmers et~al.(2023)]{dettmers2023qlora}
T.~Dettmers et~al.
\newblock QLoRA: Efficient finetuning of quantized LLMs.
\newblock In \emph{Proc.\ NeurIPS}, 2023.

\bibitem[Zhang et~al.(2023)]{zhang2023adalora}
Q.~Zhang et~al.
\newblock AdaLoRA: Adaptive budget allocation for parameter-efficient fine-tuning.
\newblock In \emph{Proc.\ ICLR}, 2023.

\bibitem[Liu et~al.(2024a)]{liu2024dora}
S.~Liu et~al.
\newblock DoRA: Weight-decomposed low-rank adaptation.
\newblock In \emph{Proc.\ ICML}, 2024.

\bibitem[Zhao et~al.(2024)]{zhao2024galore}
J.~Zhao et~al.
\newblock GaLore: Memory-efficient LLM training by gradient low-rank projection.
\newblock In \emph{Proc.\ ICML}, 2024.

\bibitem[Liu et~al.(2024b)]{liu2024loraattack}
H.~Liu et~al.
\newblock LoRA-as-an-Attack: Piercing LLM safety under the share-and-play scenario.
\newblock \emph{arXiv:2403.00108}, 2024.

\bibitem[Lermen et~al.(2023)]{lermen2023lora}
S.~Lermen et~al.
\newblock LoRA fine-tuning efficiently undoes safety training in Llama 2-Chat 70B.
\newblock \emph{arXiv:2310.20624}, 2023.

\bibitem[Huang et~al.(2024)]{huang2024composite}
H.~Huang et~al.
\newblock Composite backdoor attacks against LLMs.
\newblock In \emph{Findings of NAACL}, 2024.

\bibitem[Wei et~al.(2024)]{wei2024brittleness}
B.~Wei et~al.
\newblock Assessing the brittleness of safety alignment through pruning and low-rank modifications.
\newblock In \emph{Proc.\ ICML}, 2024.

\bibitem[Sun et~al.(2025)]{sun2025peftguard}
Z.~Sun et~al.
\newblock PEFTGuard: Detecting backdoor attacks against parameter-efficient fine-tuning.
\newblock In \emph{IEEE S\&P}, 2025.

\bibitem[Schulz et~al.(2025)]{schulz2025shadowgenes}
M.~Schulz et~al.
\newblock ShadowGenes: Computational-graph fingerprinting for model genealogy.
\newblock \emph{arXiv:2502.04321}, 2025.

\bibitem[Cohen(2024)]{cohen2024jfrog}
JFrog Security Research.
\newblock Malicious models on Hugging Face Hub. Industry report, 2024.

\bibitem[Kellas et~al.(2025)]{kellas2025pickleball}
A.~D.~Kellas, N.~Christou, W.~Jiang, P.~Li, L.~Simon, Y.~David,
V.~P.~Kemerlis, J.~C.~Davis, and J.~Yang.
\newblock PickleBall: Secure deserialization of pickle-based machine
learning models.
\newblock In \emph{Proc.\ ACM CCS}, 2025. arXiv:2508.15987.

\bibitem[Kurita et~al.(2020)]{kurita2020weight}
K.~Kurita et~al.
\newblock Weight poisoning attacks on pre-trained models.
\newblock In \emph{Proc.\ ACL}, 2020.

\bibitem[Li et~al.(2025)]{li2025backdoorllm}
Y.~Li et~al.
\newblock BackdoorLLM: A comprehensive benchmark for backdoor attacks and defenses on LLMs.
\newblock In \emph{Proc.\ NeurIPS}, 2025.

\bibitem[McMahan et~al.(2017)]{mcmahan2017fedavg}
B.~McMahan et~al.
\newblock Communication-efficient learning of deep networks from decentralized data.
\newblock In \emph{Proc.\ AISTATS}, 2017.

\bibitem[Li et~al.(2020)]{li2020fedprox}
T.~Li et~al.
\newblock Federated optimization in heterogeneous networks.
\newblock In \emph{Proc.\ MLSys}, 2020.

\bibitem[Karimireddy et~al.(2020)]{karimireddy2020scaffold}
S.~P. Karimireddy et~al.
\newblock SCAFFOLD: Stochastic controlled averaging for federated learning.
\newblock In \emph{Proc.\ ICML}, 2020.

\bibitem[Abadi et~al.(2016)]{abadi2016deep}
M.~Abadi et~al.
\newblock Deep learning with differential privacy.
\newblock In \emph{Proc.\ CCS}, 2016.

\bibitem[Mironov(2017)]{mironov2017renyi}
I.~Mironov.
\newblock R\'enyi differential privacy.
\newblock In \emph{Proc.\ CSF}, 2017.

\bibitem[Mironov et~al.(2019)]{mironov2019sampled}
I.~Mironov, K.~Talwar, and L.~Zhang.
\newblock R\'enyi differential privacy of the sampled Gaussian mechanism.
\newblock \emph{arXiv:1908.10530}, 2019.

\bibitem[Dong et~al.(2022)]{dong2022gaussian}
J.~Dong, A.~Roth, and W.~J. Su.
\newblock Gaussian differential privacy.
\newblock \emph{J. R. Stat.\ Soc.\ B}, 84(1):3--37, 2022.

\bibitem[Gopi et~al.(2021)]{gopi2021numerical}
S.~Gopi et~al.
\newblock Numerical composition of differential privacy.
\newblock In \emph{Proc.\ NeurIPS}, 2021.

\bibitem[Bonawitz et~al.(2017)]{bonawitz2017practical}
K.~Bonawitz et~al.
\newblock Practical secure aggregation for privacy-preserving machine learning.
\newblock In \emph{Proc.\ CCS}, 2017.

\bibitem[Blanchard et~al.(2017)]{blanchard2017machine}
P.~Blanchard et~al.
\newblock Machine learning with adversaries: Byzantine tolerant gradient descent.
\newblock In \emph{Proc.\ NeurIPS}, 2017.

\bibitem[Yin et~al.(2018)]{yin2018byzantine}
D.~Yin et~al.
\newblock Byzantine-robust distributed learning.
\newblock In \emph{Proc.\ ICML}, 2018.

\bibitem[Cao et~al.(2021)]{cao2021fltrust}
X.~Cao et~al.
\newblock FLTrust: Byzantine-robust federated learning by trust bootstrapping.
\newblock In \emph{Proc.\ NDSS}, 2021.

\bibitem[Anaedevha et~al.(2026a)]{anaedevha2026byzantine}
R.~N.~Anaedevha, A.~G.~Trofimov, and Y.~V.~Borodachev.
\newblock Byzantine-resilient stochastic games for federated
multi-cloud intrusion detection.
\newblock \emph{Complex \& Intelligent Systems}, Springer, 2026.
In press.
\newblock \url{https://doi.org/10.1007/s40747-026-02412-2}

\bibitem[Fan et~al.(2023)]{fan2023fatellm}
T.~Fan et~al.
\newblock FATE-LLM: A federated learning framework for large language models.
\newblock \emph{arXiv:2310.10049}, 2023.

\bibitem[Wang et~al.(2024)]{wang2024flora}
Z.~Wang et~al.
\newblock FLoRA: Federated fine-tuning large language models with heterogeneous low-rank adaptations.
\newblock In \emph{Proc.\ NeurIPS}, 2024.

\bibitem[Ye et~al.(2024)]{ye2024openfedllm}
R.~Ye, W.~Wang, J.~Chai, D.~Li, Z.~Li, Y.~Xu, Y.~Du, Y.~Wang, and S.~Chen.
\newblock OpenFedLLM: Training large language models on decentralized
private data via federated learning.
\newblock In \emph{Proc.\ ACM SIGKDD}, 2024. arXiv:2402.06954.

\bibitem[Tao et~al.(2026)]{tao2026safefedllm}
M.~Tao, Y.~Tian, W.~Tu, Y.~Yang, X.~Yang, and X.~Tang.
\newblock Safe-FedLLM: Delving into the safety of federated large
language models.
\newblock In \emph{Proc.\ ACL}, 2026. arXiv:2601.07177.

\bibitem[Kumar et~al.(2019)]{kumar2019jodie}
S.~Kumar et~al.
\newblock Predicting dynamic embedding trajectory in temporal interaction networks.
\newblock In \emph{Proc.\ KDD}, 2019.

\bibitem[Trivedi et~al.(2019)]{trivedi2019dyrep}
R.~Trivedi et~al.
\newblock DyRep: Learning representations over dynamic graphs.
\newblock In \emph{Proc.\ ICLR}, 2019.

\bibitem[Xu et~al.(2020)]{xu2020tgat}
D.~Xu et~al.
\newblock Inductive representation learning on temporal graphs.
\newblock In \emph{Proc.\ ICLR}, 2020.

\bibitem[Rossi et~al.(2020)]{rossi2020tgn}
E.~Rossi et~al.
\newblock Temporal graph networks for deep learning on dynamic graphs.
\newblock In \emph{Proc.\ ICML Workshop on Graph Representation Learning}, 2020.

\bibitem[Yu et~al.(2023)]{yu2023dygformer}
L.~Yu et~al.
\newblock Towards better dynamic graph learning: New architecture and unified library.
\newblock In \emph{Proc.\ NeurIPS}, 2023.

\bibitem[Anaedevha and Trofimov(2026)]{anaedevha2026multiscale}
R.~N.~Anaedevha and A.~G.~Trofimov.
\newblock Multi-scale temporal graph neural network with stochastic
dynamics for cloud and industrial network intrusion detection.
\newblock In \emph{Proc.\ IEEE EDM}, pp.\ 1--8, 2026.
\newblock \url{https://doi.org/10.1109/EDM69524.2026.11632205}

\bibitem[Hu et~al.(2020)]{hu2020hgt}
Z.~Hu et~al.
\newblock Heterogeneous graph transformer.
\newblock In \emph{Proc.\ WWW}, 2020.

\bibitem[He et~al.(2021)]{he2021fedgraphnn}
C.~He et~al.
\newblock FedGraphNN: A federated learning benchmark system for graph neural networks.
\newblock In \emph{Proc.\ ICLR Workshop}, 2021.

\bibitem[Sajadmanesh et~al.(2023)]{sajadmanesh2023gap}
S.~Sajadmanesh et~al.
\newblock GAP: Differentially-private graph neural networks with aggregation perturbation.
\newblock In \emph{USENIX Security}, 2023.

\bibitem[Rahman et~al.(2025)]{rahman2025hugginggraph}
M.~S.~Rahman, P.~Gao, and Y.~Ji.
\newblock HuggingGraph: Understanding the supply chain of the LLM
ecosystem.
\newblock In \emph{Proc.\ ACM CIKM}, 2025. arXiv:2507.14240.

\bibitem[Horwitz et~al.(2025)]{horwitz2025atlas}
E.~Horwitz, N.~Kurer, J.~Kahana, L.~Amar, and Y.~Hoshen.
\newblock Charting and navigating Hugging Face's model atlas.
\newblock \emph{arXiv:2503.10633}, 2025.

\bibitem[Lecuyer et~al.(2019)]{lecuyer2019certified}
M.~Lecuyer et~al.
\newblock Certified robustness to adversarial examples with differential privacy.
\newblock In \emph{IEEE S\&P}, 2019.

\bibitem[Cohen et~al.(2019)]{cohen2019certified}
J.~Cohen, E.~Rosenfeld, and J.~Z. Kolter.
\newblock Certified adversarial robustness via randomized smoothing.
\newblock In \emph{Proc.\ ICML}, 2019.

\bibitem[Ma et~al.(2019)]{ma2019data}
Y.~Ma et~al.
\newblock Data poisoning against differentially-private learners.
\newblock In \emph{Proc.\ IJCAI}, 2019.

\bibitem[Anaedevha et~al.(2026b)]{anaedevha2026mambashield}
R.~N.~Anaedevha, A.~G.~Trofimov, and Y.~V.~Borodachev.
\newblock MambaShield: Selective state-space models for
poisoning-resilient sequential modeling.
\newblock \emph{Expert Systems with Applications}, p.~131175,
Elsevier, 2026.
\newblock \url{https://doi.org/10.1016/j.eswa.2026.131175}

\bibitem[Xie et~al.(2021)]{xie2021crfl}
C.~Xie et~al.
\newblock CRFL: Certifiably robust federated learning against backdoor attacks.
\newblock In \emph{Proc.\ ICML}, 2021.

\bibitem[Cao et~al.(2022)]{cao2022flcert}
X.~Cao et~al.
\newblock FLCert: Provably secure federated learning against poisoning attacks.
\newblock In \emph{IEEE TIFS}, 2022.

\bibitem[Wang et~al.(2021)]{wang2021certgnn}
B.~Wang et~al.
\newblock Certified robustness of graph neural networks against adversarial structural perturbation.
\newblock In \emph{Proc.\ KDD}, 2021.

\bibitem[Scholten et~al.(2022)]{scholten2022randomized}
Y.~Scholten et~al.
\newblock Randomized message-interception smoothing.
\newblock In \emph{Proc.\ NeurIPS}, 2022.

\bibitem[Beutel et~al.(2020)]{beutel2020flower}
D.~J. Beutel et~al.
\newblock Flower: A friendly federated learning research framework.
\newblock \emph{arXiv:2007.14390}, 2020.

\bibitem[Loshchilov and Hutter(2019)]{loshchilov2019adamw}
I.~Loshchilov and F.~Hutter.
\newblock Decoupled weight decay regularization.
\newblock In \emph{Proc.\ ICLR}, 2019.

\bibitem[Anaedevha et~al.(2026c)]{anaedevha2026gp}
R.~N.~Anaedevha, A.~G.~Trofimov, and Y.~V.~Borodachev.
\newblock Uncertainty-calibrated hierarchical Gaussian processes for
intrusion detection with multi-scale temporal modeling.
\newblock \emph{Neurocomputing}, 677:133105, Elsevier, 2026.
\newblock \url{https://doi.org/10.1016/j.neucom.2026.133105}

\bibitem[Sharafaldin et~al.(2018)]{sharafaldin2018toward}
I.~Sharafaldin et~al.
\newblock Toward generating a new intrusion detection dataset.
\newblock In \emph{Proc.\ ICISSP}, 2018.

\bibitem[Ferrag et~al.(2022)]{ferrag2022edge}
M.~A. Ferrag et~al.
\newblock Edge-IIoTset.
\newblock \emph{IEEE Access}, 10:40281--40306, 2022.

\bibitem[Moustafa and Slay(2015)]{moustafa2015unsw}
N.~Moustafa and J.~Slay.
\newblock UNSW-NB15.
\newblock In \emph{Proc.\ MilCIS}, 2015.

\bibitem[Moustafa(2021)]{moustafa2021ton}
N.~Moustafa.
\newblock TON-IoT telemetry dataset.
\newblock \emph{IEEE Access}, 9:55141--55167, 2021.

\end{thebibliography}

% ============================================================
%  APPENDICES
% ============================================================
\appendix

\section{Full Proofs}
\label{app:proofs}

\subsection{Proof of Theorem~\ref{thm:sensitivity}}

\paragraph{Two-sided clipping bound.}
The client-update clipping map
$T_S(\mathbf{g}):=\mathbf{g}\cdot\min(1,S/\norm{\mathbf{g}}_2)$
satisfies $\norm{T_S(\mathbf{g})}_2\le S$ for every
$\mathbf{g}\in\R^d$. Hence for any two datasets $\Dcal_i$ and
$\Dcal_i'$,
$\norm{\bar{\mathbf{g}}_i(\Dcal_i)-\bar{\mathbf{g}}_i(\Dcal_i')}_2
\le \norm{\bar{\mathbf{g}}_i(\Dcal_i)}_2 +
\norm{\bar{\mathbf{g}}_i(\Dcal_i')}_2 \le 2S$, giving
Eq.~\eqref{eq:sensitivity}. This is the sensitivity we use for the
Gaussian mechanism calibration and does not depend on any
assumption on the GNN architecture.

\paragraph{Graph-structural refinement.}
When the sensitivity to a single node-edge change propagates through
$L$ HGT layers rather than saturating clipping, we can bound
$\norm{\Delta \bar{\mathbf{g}}_i}_2$ more tightly. By induction on
$L$: at the $\ell$-th HGT aggregation (Eq.~\ref{eq:hgt}),
$\mathbf{h}_v^{(\ell)} = \phi^{\mathrm{nl}}(\sum_{r,u} \alpha_{u,v,r}
\mathbf{W}_R^r \mathbf{h}_u^{(\ell-1)})$, and differing in a
single node-edge changes at most one term in the inner sum. By
Cauchy--Schwarz and the relation-wise spectral bound
$\norm{\mathbf{W}_R^r}_2 \le B_W$,
\begin{equation}
\norm{\Delta \mathbf{h}_v^{(\ell)}}_2 \le \rho \cdot B_W \cdot d_{\max} \cdot \norm{\Delta \mathbf{h}_v^{(\ell-1)}}_2
\end{equation}
whenever the affected node lies in the temporal neighbourhood.
Iterating $L$ layers yields the $(\rho B_W d_{\max})^L$ factor.
The $|\Rcal|^{1/2}$ factor comes from Cauchy--Schwarz over the
relation-type sum. Post-clipping to $S$ then gives the refined
Eq.~\eqref{eq:sensitivity-refined}, and by taking the tighter of
the two bounds we obtain
$\Delta_2\le\min\!\bigl(2S,2S(\rho B_W d_{\max})^L |\Rcal|^{1/2}\bigr)$.
For our deployment scale
$(\rho B_W d_{\max})^L\,|\Rcal|^{1/2}\gg 1$, so the operational
sensitivity is $2S$.
\qed

\subsection{Proof of Theorem~\ref{thm:certified}}

The argument composes three established results, itemised below
with the specific hypothesis each contributes.

\paragraph{Step 1: PixelDP stability of the smoothed classifier.}
By Lemma~1 of \citet{lecuyer2019certified}, if $f$ is an
$(\epsilon_r,\delta)$-DP scoring function of its input
$\mathbf{x}$, then for any input perturbation $\Delta$ with
bounded $\ell_2$ norm,
$\Prob[\argmax f(\mathbf{x}) \neq \argmax f(\mathbf{x} +
\Delta)] \leq e^{\epsilon_r} \cdot p_{\Delta}(\mathbf{x}) + \delta$,
where $p_{\Delta}$ is the prediction-flip probability under the
noiseless mechanism. In our setting the random variable being
smoothed is the aggregated global model $\mathbf{w}_T\in\R^d$
(not the input adapter), and the $(\epsilon_r,\delta)$-DP mechanism
is the single Gaussian smoothing draw applied to $\mathbf{w}_T$ at
verification time.

\paragraph{Step 2: from parameter perturbation to client coalition.}
A single malicious client $i$ operating under
Definition~\ref{def:attack} submits
$\norm{\tilde{\mathbf{g}}_i-\mathbf{g}_i^{\mathrm{honest}}}_2\le
\epsilon_{\mathrm{wt}}$ after SecAgg. FLTrust weighting
$\beta_i\in[0,1]$ can only reduce the effective attacker
contribution to the aggregate, so any bound derived under
unweighted aggregation is a fortiori valid under FLTrust
weighting. Under unweighted averaging, $k$ colluding clients induce
a worst-case parameter-space perturbation of the aggregate of size
at most $\eta\cdot k\cdot\epsilon_{\mathrm{wt}}/m$ per round by
triangle inequality; here $\eta$ is the learning rate and $m$ is
the number of clients sampled per round. Summed over the training
horizon the total attacker-induced parameter perturbation is thus
bounded linearly in $k$.

\paragraph{Step 3: CRFL federated smoothing.}
By Theorem~1 of \citet{xie2021crfl}, extending randomized
smoothing~\cite{cohen2019certified}, federated parameter clipping
plus a Gaussian smoothing draw on the aggregated model certifies
a class-margin radius (in parameter-perturbation space) via a
Neyman--Pearson hypothesis-testing argument. Translated into a
client coalition of size $N$, this reads
\begin{equation}
k^* = \left\lfloor \frac{N (\Phi^{-1}(\hat{p}_1) - \Phi^{-1}(\hat{p}_2))}{2 e^{\epsilon_r}} \right\rfloor,
\end{equation}
where $\epsilon_r$ is the $(\epsilon_r,\delta)$-DP budget of the
single smoothing draw applied at verification time. The factor
$2 e^{\epsilon_r}$ combines the PixelDP per-query stability
multiplier $e^{\epsilon_r}$ (Step~1) with the CRFL
class-margin denominator $2$ (two half-Gaussian tails). The
PixelDP factor is therefore instantiated at $\epsilon_r$, not at
the cumulative training-time budget $\epsilon_T$.

\paragraph{Step 4: R\'enyi-DP composition (marketplace privacy only).}
By the standard R\'enyi-DP composition
theorem~\cite{mironov2017renyi}, the Gaussian mechanism satisfies
$(\alpha,\alpha/(2\sigma^2))$-RDP, and for the
\emph{Poisson-subsampled} Gaussian mechanism at rate $q$ the
per-round R\'enyi divergence is amplified to
$O(q^2\alpha/\sigma^2)$~\cite{mironov2019sampled}; composing $T$
rounds and converting to $(\epsilon_T,\delta)$-DP gives
\begin{equation}
\epsilon_T \leq \min_{\alpha>1}\;\frac{2 T q^2 \alpha}{\sigma^2} + \frac{\ln(1/\delta)}{\alpha - 1}.
\end{equation}
In practice we do not use this closed form: at $\sigma=1.1$,
$q=0.2$, $T=100$, $\delta=10^{-5}$, the Opacus PRV
accountant~\cite{gopi2021numerical} returns
$\epsilon_T\approx 13.6$, and we report this numerical value as
the marketplace-privacy guarantee. This step does \emph{not} enter
the certificate: the certificate is set by the per-query smoothing
budget $\epsilon_r$ of Step~3, not by the cumulative training-time
$\epsilon_T$ of Step~4.

Substituting the per-query smoothing budget $\epsilon_r$ into the
CRFL bound yields Eq.~\eqref{eq:kstar}. The failure probability
$T\delta$ accumulates by union bound over rounds if the smoothing
is re-drawn per round; a single-draw certificate uses the tighter
$\delta$. The graph-structural extension to heterogeneous edge-set
perturbations follows from
\citet{wang2021certgnn,scholten2022randomized}.
\qed

\section{Additional Analyses}
\label{app:additional}

\paragraph{Spectral-only baseline.}
Because backdoored adapters carry a prominent rank-1 signature
(Figure~\ref{fig:spectrum}), we test whether that signal alone
suffices. A non-graph, non-federated SVD-only detector thresholding
$\sigma_1/\sigma_2$ reaches $90.2\%$ macro-F1 on single-trigger
rank-1 backdoors but only $71.4\%$ on composite-trigger,
weight-poisoning, and trigger-free merging attacks (the
logistic-regression variant on the full spectral signature is
similar: $91.0\%/73.1\%$), and offers no certificate. \modelname's
advantage is therefore concentrated on the hard, context-dependent
attacks and on the certified radius, not on the easy rank-1 cases
where a cheap spectral screen suffices.

\paragraph{Leakage-controlled splits.}
Table~\ref{tab:splits} reports \modelname's macro-F1 under four
increasingly strict test partitions of \benchname. Under a random
adapter split (the standard split) \modelname\ reaches $96.4\%$;
under a lineage-component-disjoint split (the strictest split that
still preserves attack coverage) it retains $95.1\%$; under a
base-model-disjoint split it drops to $93.7\%$; under an
attack-family-held-out split (each attack type in turn) it retains
$92.3\%$ averaged over held-out families. The main-paper number
$96.4\%$ is reported at the lineage-component-disjoint split, so
it is not an artefact of family memorisation.

\begin{table}[!h]
\centering
\caption{Macro-F1 of \modelname\ under increasingly strict test
partitions of \benchname.}
\label{tab:splits}
\small
\begin{tabular}{@{}lc@{}}
\toprule
Split                                   & Macro-F1 (\%) \\
\midrule
Random adapter                          & 96.4 \\
Lineage-component-disjoint (main)       & 95.1 \\
Base-model-disjoint                     & 93.7 \\
Attack-family-held-out (avg)            & 92.3 \\
\bottomrule
\end{tabular}
\end{table}

\paragraph{Security-operating-point metrics.}
Because macro-F1 on the balanced $50/50$ split can under-report
the false-positive burden a marketplace would actually see, we
also report precision, recall, false-positive rate (FPR), and
detection@FPR$=1\%$ at the paper's default operating point of
$\hat{p}_{\mathrm{mal}}\ge 0.6$. Values (mean over three seeds):
precision $= 94.4\%$, recall $= 96.9\%$, FPR $= 5.9\%$, PR-AUC
$= 0.972$, detection@FPR$=1\%$ $= 82.4\%$. On the imbalanced
production distribution (estimated $\sim 0.4\%$ malicious base
rate in prior industry reports) the operator would tune to a
low-FPR operating point; at $\hat{p}_{\mathrm{mal}}\ge 0.9$ we
observe FPR $= 0.7\%$ with detection rate $84.1\%$.

\paragraph{Robustness of the operating envelope.}
At a stronger-privacy point (cumulative $\epsilon_T = 1.0$) macro-F1
is $94.1\%$, a $2.3$\,pp drop from the operational $\epsilon_T
\approx 13.6$ point --- graceful degradation, giving a
privacy--utility curve rather than a single point. Under random
removal of HuggingGraph lineage edges, macro-F1 falls from
$96.4\%$ to $95.6\%/93.8\%/91.1\%$ at $10\%/30\%/50\%$ removal, so
a complete lineage graph is not required. A two-sample
Kolmogorov--Smirnov test on the top-1 singular-value distributions
confirms the $2.43\times$ benign/backdoored separation is
significant ($D=0.71$, $p<10^{-3}$). Latency is $0.7$\,s per
adapter on an A100~80\,GB (including neighbour retrieval and
feature extraction), rising to $\sim\!1.6$\,s on an A10G and
$\sim\!2.9$\,s on a T4; the ``sub-second'' claim is thus specific
to A100-class hardware.

\section{IDS Transfer Results}
\label{app:ids}

As a cross-domain sanity check for the detector \emph{backbone}
beyond the adapter-integrity setting (i.e., emphatically not part
of the LoRA-integrity contribution), we evaluate on four
network intrusion-detection datasets --- CIC-IDS2017~\cite{sharafaldin2018toward},
Edge-IIoTset~\cite{ferrag2022edge}, UNSW-NB15~\cite{moustafa2015unsw},
and TON-IoT~\cite{moustafa2021ton}, following the evaluation
protocol of prior uncertainty-aware intrusion-detection
work~\cite{anaedevha2026gp}. Averaged over the four,
\modelname\ attains $96.1\%$ accuracy, $95.8\%$ macro-F1, and
expected calibration error $0.034$, versus $91.8\%/0.072$ for
FedAvg-MLP, $93.6\%/0.054$ for FedGraphNN-HGT, and
$94.1\%/0.049$ for Krum-DyGFormer. These results indicate the
federated dynamic-graph backbone transfers to non-adapter
anomaly-detection settings, but they are not part of the core
LoRA-integrity contribution and are reported here only for
completeness.

\section{Benchmark Provenance}
\label{app:provenance}

Table~\ref{tab:provenance} audits which components of \benchname\
are drawn from real released artefacts versus synthesised by the
benchmark construction pipeline.

\begin{table}[!h]
\centering
\caption{Provenance audit of \benchname: real (drawn from public
released artefacts) versus synthetic (produced by the benchmark
construction pipeline). The synthetic components are all seeded
and their generation code is released so the audit is
reproducible.}
\label{tab:provenance}
\footnotesize
\begin{tabular}{@{}lccl@{}}
\toprule
Component                     & Real & Synth. & Source \\
\midrule
Base model families           & \checkmark &            & HF Hub \\
LoRA adapter weights          & \checkmark & \checkmark & PADBench + this work \\
Attack recipes                & \checkmark &            & BackdoorLLM, PADBench \\
Model cards / README text     & \checkmark & \checkmark & HF + templated \\
Contributor identities        &            & \checkmark & this work \\
Derivation (lineage) edges    & \checkmark & \checkmark & HuggingGraph + this work \\
Deployment edges              &            & \checkmark & this work \\
Citation edges                & \checkmark & \checkmark & PADBench + this work \\
Temporal timestamps           &            & \checkmark & this work \\
Download / usage statistics   & \checkmark & \checkmark & HF + templated \\
\bottomrule
\end{tabular}
\end{table}

\benchname\ is released with SHA-256 checksums for every artefact,
a manifest listing which fields per adapter are real versus
synthesised, and the generation seeds so the synthetic components
are reproducible bit-for-bit.

\section{Reproducibility}
\label{app:repro}

All results are means over three random seeds $\{42, 137, 2026\}$.
Total compute: approximately $960$ A100-GPU-hours across
hyperparameter sweeps. The codebase, the \benchname\ benchmark,
and pretrained checkpoints are released at the repository below.
Hyperparameters: DyGFormer encoder $3$ layers, $H = 8$ attention
heads, $d_v = 256$. HGT relation-specific matrices initialised
with Xavier-uniform. Training: AdamW with learning rate $5 \times
10^{-4}$, weight decay $10^{-4}$, $T = 100$ rounds, batch size
$32$ per client. Client-update-level DP-SGD: $S = 1.0$, $\sigma =
1.1$, Poisson sampling rate $q = 0.2$, $\delta = 10^{-5}$. FLTrust
root set $|R_0| = 200$ vetted clean adapters. Per-query smoothing
budget for the certificate: $\epsilon_r = 0.5$.

\paragraph{DP accounting configuration.} The reported
$\epsilon_T \approx 13.6$ at $T=100$ rounds is computed with the
Opacus PRV accountant~\cite{gopi2021numerical} under
\emph{Poisson} subsampling at sampling rate $q = 0.2$, noise
multiplier $\sigma = 1.1$, client-update clipping $S = 1.0$,
$\delta = 10^{-5}$. As a cross-check, the Opacus R\'enyi
accountant returns $\epsilon_T \approx 13.56$ at the same
parameters --- consistent with PRV within accountant looseness.
The manuscript's earlier draft misreported $\epsilon_T\approx 5.0$
for these parameters; that value would require $\sigma\approx
2.15$ at $(q,T,\delta)=(0.2,100,10^{-5})$ under the same PRV
accountant. We report the honest PRV output for the sigma
actually used in the released code and checkpoints. The per-query
certificate smoothing budget $\epsilon_r=0.5$ is a separate design
parameter of the inference-time smoothing draw and is not
composed into $\epsilon_T$; it enters only Theorem~\ref{thm:certified}.
The exact configuration is::
\begin{quote}
\small
\begin{verbatim}
accountant       = PRV (Opacus)
sample_rate q    = 0.20 (Poisson)
sigma            = 1.10
clip S           = 1.00
T                = 100 rounds
delta            = 1e-5
result eps_T     = 13.6 (PRV), 13.56 (RDP)
\end{verbatim}
\end{quote}

\paragraph{Training-time overhead of the privacy stack.}
Because \modelname\ applies DP clipping at the
\emph{client-update} level (Algorithm~\ref{alg:fedloraguard},
line~5) rather than to per-example gradients, it does not require
the per-sample gradient materialisation or microbatching that
dominates the cost of example-level DP-SGD: each round adds one
$O(d)$ norm-and-rescale and one $O(d)$ Gaussian draw per client,
where $d$ is the encoder parameter count. Secure aggregation
contributes a pairwise PRG mask that is likewise $O(d)$ per
client, plus the mask-exchange round of the Bonawitz
protocol~\cite{bonawitz2017practical}, whose cost is
communication- rather than compute-bound and is amortised over the
local epochs of a round. We instrument the released runtime to
measure this overhead directly. On identical CPU hardware, over
three independent runs with $20$--$60$ measurement rounds each
after $3$--$5$ warm-up rounds and identical client work (six
clients, $H{=}4$ HGT heads, $d_v=96$, batch size $32$), the
pooled median round wall-clock is $1.89$\,s under non-private
FedAvg and $1.96$\,s with the full DP + secure-aggregation stack,
a pooled median overhead of $+3.4\%$ (across the three runs
individually the full-stack overhead is $-1.1\%$, $+10.6\%$, and
$+3.4\%$; per-round overhead of the DP step alone is dominated by
CPU-jitter at this model scale, so we round the reported figure
to $\approx 4\%$). Because CPU-hosted federation puts most of the
per-round cost in local GNN training, the DP + SecAgg mask draws
account for a small fraction of round time; on
compute-accelerated (A100-class) hardware the same mask draws
would represent a larger relative share and the ratio would be
higher. Convergence in rounds is unchanged ($T = 100$); the
certificate and privacy guarantee are therefore obtained at a
modest training-time cost.
At verification time, no additional privacy training step is
required beyond the forward pass and the single per-query
smoothing draw described in Theorem~\ref{thm:certified}; the
end-to-end per-adapter verification budget on A100 is $0.7$\,s
(Appendix~\ref{app:additional}). We ship the round-time
instrumentation in the released runtime so this overhead can be
re-measured on other hardware.

\paragraph{Availability.} The codebase, the \benchname\ benchmark,
pinned dependency hashes, and pretrained checkpoints are released
at \url{https://github.com/rogerpanel/CV/tree/main/FedLoRAGuard}
(reference implementation and reproducibility artefact) and
\url{https://github.com/rogerpanel/FedLoRAGuard-Models}
(manuscript source). Dataset sources and licences are cited at
first use and listed in the Ethics Statement.

\end{document}
